% DERIVED from formal_layer.json — read-only, do not edit.
\begin{definitionv}{synth_11}[Complete observed-record space]
For each natural number \(d\), define the complete observed-record space by
\[
\operatorname{Obs}(d)\coloneqq [d]\times\{0,1\}\times\{0,1\}.
\]
An element of \(\operatorname{Obs}(d)\) is a triple \((X,A,Y)\), where \(X\in[d]\) is the finite-alphabet covariate, \(A\in\{0,1\}\) is the binary treatment indicator, and \(Y\in\{0,1\}\) is the binary observed outcome.
\end{definitionv}

\begin{definitionv}{synth_16}[Covariate cell mass]
For an observed-data law \(P\) on \([d]\times\{0,1\}^2\) and a covariate value \(x\in[d]\), define the covariate cell mass
\[
p_x
=
\sum_{a\in\{0,1\}}\sum_{y\in\{0,1\}}
P(X=x,A=a,Y=y).
\]
\end{definitionv}

\begin{definitionv}{synth_10}[Cell treatment propensity]
For an observed-data law \(P\) on \([d]\times\{0,1\}^2\) and a covariate cell \(x\in[d]\), define the treatment propensity \(e_x\) by
\[
e_x
:=
\frac{\sum_{y\in\{0,1\}} P(X=x,A=1,Y=y)}
{\sum_{a\in\{0,1\}}\sum_{y\in\{0,1\}} P(X=x,A=a,Y=y)}.
\]
The ratio is interpreted with the usual totalized real division convention, so it is \(0\) when the denominator is \(0\).
\end{definitionv}

\begin{assumptionv}{ass:overlap}[Occupied-cell overlap]
For every covariate cell \(x\in[d]\) with cell mass \(p_x>0\), the propensity score satisfies
\[
\epsilon \le e_x \le 1-\epsilon .
\]
\end{assumptionv}

\begin{assumptionv}{ass:labeled-iid}[Labeled sampling]
The labeled sample is distributed as
\[
\mathcal L_n \sim P^{\otimes n}.
\]
\end{assumptionv}

\begin{definitionv}{synth_14}[Treatment-covariate marginal]
For any \(d\) and any observed-data law \(P\) on \([d]\times\{0,1\}^2\), the treatment-covariate marginal \(P_{XA}\) is the probability mass function on \([d]\times\{0,1\}\) obtained by pushing \(P\) forward under the projection \((x,a,y)\mapsto (x,a)\):
\[
P_{XA}
=
P\circ\bigl((x,a,y)\mapsto (x,a)\bigr)^{-1}.
\]
\end{definitionv}

\begin{assumptionv}{ass:auxiliary-iid}[Auxiliary sampling]
The auxiliary sample is distributed as
\[
\mathcal U_m \sim P_{XA}^{\otimes m}.
\]
\end{assumptionv}

\begin{assumptionv}{ass:sample-independence}[Sample independence]
The labeled and auxiliary samples are independent:
\[
\mathcal L_n \perp\!\!\!\perp \mathcal U_m .
\]
\end{assumptionv}

\begin{assumptionv}{ass:consistency}[Potential-outcome consistency]
The observed binary outcome equals the potential outcome under the realized treatment:
\[
Y = Y(A)
\]
almost surely.
\end{assumptionv}

\begin{assumptionv}{ass:conditional-exchangeability}[Conditional exchangeability]
The potential outcomes are conditionally independent of treatment given the finite covariate:
\[
(Y(0),Y(1)) \perp\!\!\!\perp A \mid X .
\]
\end{assumptionv}

\begin{definitionv}{synth_13}[Probability laws on a measurable space]
For a measurable space \(\mathcal Z\), let \(\mathsf{Prob}(\mathcal Z)\) denote the set of all probability measures on \(\mathcal Z\). For the finite spaces used below, including \([d]\times\{0,1\}^2\) and \([d]\times\{0,1\}\), the measurable structure is the full power set.
\end{definitionv}

\begin{definitionv}{def:model-class}[Model class \(\mathcal M_{d,\epsilon}\)]
For an overlap constant \(\epsilon\) and alphabet size \(d\), define
\[
\mathcal M_{d,\epsilon}
:=
\Bigl\{
P\in\mathsf{Prob}\bigl([d]\times\{0,1\}^2\bigr):
d\ge2,\ 0<\epsilon<1/2,\ \text{\(P\) satisfies \cref{obj:ass:overlap}}
\Bigr\}.
\]
\end{definitionv}

\begin{definitionv}{synth_17}[Outcome regression]
For a law \(P\) on \([d]\times\{0,1\}^2\), a treatment arm \(a\in\{0,1\}\), and a covariate cell \(x\in[d]\), define
\[
\mu_{ax}(P)
:=
\frac{P(X=x,A=a,Y=1)}
{\sum_{y\in\{0,1\}} P(X=x,A=a,Y=y)} .
\]
The quotient is interpreted as total real division, so \(\mu_{ax}(P)=0\) when the treatment--covariate cell mass \(\sum_{y\in\{0,1\}}P(X=x,A=a,Y=y)\) is zero. In particular, \(\mu_{1x}(P)\) and \(\mu_{0x}(P)\) are obtained by taking \(a=1\) and \(a=0\), respectively.
\end{definitionv}

\begin{definitionv}{def:ate-functional}[Average treatment effect functional]
For a finite alphabet size \(d\) and an observed-data law \(P\) on \([d]\times\{0,1\}^2\), define the average treatment effect functional by
\[
\tau(P):=\sum_{x=1}^{d}p_x(P)\bigl(\mu_{1x}(P)-\mu_{0x}(P)\bigr).
\]
Here \(p_x(P)\) is the covariate cell mass at \(x\), and \(\mu_{ax}(P)\) is the binary outcome regression in treatment arm \(a\) and covariate cell \(x\), totalized by division in the arm cell.
\end{definitionv}

\begin{definitionv}{def:annotation-experiment}[Annotation experiment \(\mathcal E_{n,m,d,\epsilon}\)]
Define the annotation experiment by
\[
\mathcal E_{n,m,d,\epsilon}
:=
\Bigl\{
P^{\otimes n}\otimes P_{XA}^{\otimes m}:
P\in\mathcal M_{d,\epsilon},\ n\ge1,\ m\ge0,\ d\ge2
\Bigr\}.
\]
The product law corresponds to \(n\) independent complete outcome-labeled records and \(m\) independent additional outcome-unlabeled treatment-covariate records, with independence between the two samples.
\end{definitionv}

\begin{definitionv}{synth_18}[Joint sample space]
For natural numbers \(n,m,d\), the joint sample space is
\[
\mathrm{Sample}(n,m,d)
=
\bigl([n]\to [d]\times\{0,1\}\times\{0,1\}\bigr)
\times
\bigl([m]\to [d]\times\{0,1\}\bigr).
\]
Equivalently, an element consists of \(n\) complete records \((X,A,Y)\) and \(m\) outcome-unlabeled records \((X,A)\).
\end{definitionv}

\begin{definitionv}{def:minimax-risk}[Minimax risk \(R_\epsilon(n,m,d)\)]
Define the minimax mean-squared risk by
\[
R_\epsilon(n,m,d)
:=
\inf_{T\in\mathcal T_{n,m,d}}
\sup_{P\in\mathcal M_{d,\epsilon}}
E_{P^{\otimes n}\otimes P_{XA}^{\otimes m}}
\!\left[\bigl(T-\tau(P)\bigr)^2\right],
\]
where \(\mathcal T_{n,m,d}\) is the class of estimators measurable with respect to the labeled sample and auxiliary sample.
\end{definitionv}

\begin{definitionv}{def:frontier-rate-handle}[Frontier rate \(r_\epsilon(n,m,d)\)]
Put
\[
N=n+m,
\qquad
\ell_n=\log(en).
\]
The unequal-information annotation-frontier rate is
\[
r_\epsilon(n,m,d)
:=
\min\!\left\{
1,\frac1n+\frac{d^2}{N^2\ell_n^2}
\right\}.
\]
\end{definitionv}

\begin{definitionv}{synth_7}[Finite calibration scales for the mixed estimator]
Set
\[
A_\star=7056,\qquad c^\circ=\frac1{256}.
\]
For \(0<\epsilon<1/2\), let \(j_\epsilon\) be the least integer \(j\ge1\) such that \(2^{-j}\le\epsilon\), and define the dyadic overlap certificate
\[
\bar\epsilon(\epsilon)=2^{-j_\epsilon}.
\]
Then \(0<\bar\epsilon(\epsilon)\le\epsilon<2\bar\epsilon(\epsilon)\). Let \(H_\epsilon\) be the least positive integer \(H\) satisfying
\[
H\ge4,\qquad \frac H4\ge15+\frac{24}{c^\circ},\qquad
Hc^\circ\bar\epsilon(\epsilon)\ge672,\qquad
Hc^\circ\bar\epsilon(\epsilon)\ge24.
\]
For \(n\ge1\), put
\[
h_n=1+\lceil\log_2 n\rceil,\qquad
L(n)=\max\{2,\lfloor c^\circ h_n\rfloor\}.
\]
For \(n\ge24\) and \(m\ge0\), define
\[
M_0=\lfloor n/3\rfloor,\qquad
n_p=\left\lfloor\frac{n-M_0}{2}\right\rfloor,\qquad
n_f=n-M_0-n_p,
\]
\[
m_p=\lfloor m/2\rfloor,\qquad m_f=m-m_p,\qquad
M_p=n_p+m_p,\qquad M_f=n_f+m_f,
\]
and
\[
t_p=\frac{M_p}{8},\qquad t=\frac{M_f}{8}.
\]
The light-cell scale and pilot-count threshold are
\[
B(n,m,\epsilon)=\frac{H_\epsilon L(n)}{\min(t_p,t)},\qquad
k_0(n,m,\epsilon)=\left\lfloor\frac{t_p B(n,m,\epsilon)}4\right\rfloor .
\]
For \(n<24\), the calibration predicate using these quantities is declared false.
\end{definitionv}

\begin{definitionv}{synth_12}[Induced law of the treatment--covariate marginal]
Let \(b\in\{0,1\}\) and let \(\Pi_b\) be a probability measure on \(\mathcal M_{d,\epsilon}\). If \(P\sim\Pi_b\), define
\[
\mathcal L_{\Pi_b}(P_{XA})
:=
(P\mapsto P_{XA})_{\#}\Pi_b .
\]
Thus, for every measurable set \(B\subseteq\mathsf{Prob}([d]\times\{0,1\})\),
\[
\mathcal L_{\Pi_b}(P_{XA})(B)
=
\Pi_b\{P\in\mathcal M_{d,\epsilon}:P_{XA}\in B\}.
\]
Equivalently, \(\mathcal L_{\Pi_b}(P_{XA})\) is the distribution of the full \(2d\)-cell treatment--covariate marginal table induced when the complete-data law \(P\) is drawn from \(\Pi_b\).
\end{definitionv}

\begin{definitionv}{def:known-marginal-risk}[Known-marginal minimax risk]
For nonnegative integers \(n\) and \(d\) and a real overlap parameter \(\epsilon\), the known-marginal minimax risk is
\[
R_\epsilon^{\mathrm{known}}(n,d)
:=
\inf_{T\in\mathcal T^{\mathrm{known}}_{n,d}}
\sup_{P\in\mathcal M_{d,\epsilon}}
\int
\bigl(T(\mathcal L_n,P_{XA})-\tau(P)\bigr)^2
\,dP^{\otimes n}.
\]
Here \(\mathcal M_{d,\epsilon}\) is the model class in \cref{obj:def:model-class}, \(\tau(P)\) is the average treatment effect functional in \cref{obj:def:ate-functional}, and \(\mathcal T^{\mathrm{known}}_{n,d}\) is the class of measurable estimators \(T\) based on the labeled sample together with the exact treatment-covariate table \(P_{XA}\), with \(T(\mathcal L_n,P_{XA})\in[-1,1]\) for every input.
\end{definitionv}

\begin{theoremv}{thm:known-marginal-boundary}[Known marginal boundary]
Suppose the overlap constant \(\epsilon\) satisfies \(0<\epsilon<1/2\). For the known-marginal minimax risk \(R_\epsilon^{\mathrm{known}}(n,d)\) defined in \cref{obj:def:known-marginal-risk}, there exist real constants \(c_\epsilon\) and \(C_\epsilon\) with \(0<c_\epsilon\le C_\epsilon\) such that, for every \(n,d\in\mathbb N\) with \(n\ge1\) and \(d\ge2\),
\[
\frac{c_\epsilon}{n}
\le
R_\epsilon^{\mathrm{known}}(n,d)
\le
\frac{C_\epsilon}{n}.
\]
\end{theoremv}

\begin{definitionv}{synth_8}[Binary length \(b_n\)]
\begin{itemize}
\item \textbf{(Context.)} Work under \cref{obj:synth_1,obj:synth_2,obj:synth_3,obj:synth_4,obj:synth_5,obj:synth_6,obj:synth_7,obj:synth_9,obj:synth_10,obj:synth_11,obj:synth_12,obj:synth_13,obj:synth_14,obj:synth_16,obj:synth_17,obj:synth_18,obj:synth_19,obj:synth_20,obj:synth_21,obj:synth_22,obj:ass:overlap,obj:ass:labeled-iid,obj:ass:auxiliary-iid,obj:ass:sample-independence,obj:ass:consistency,obj:ass:conditional-exchangeability,obj:def:model-class,obj:def:ate-functional,obj:def:annotation-experiment,obj:def:minimax-risk,obj:def:frontier-rate-handle,obj:def:mixed-estimator-handle,obj:def:common-marginal-prior-handle,obj:def:finite-small-instance,obj:def:known-marginal-experiment,obj:def:known-marginal-risk,obj:thm:known-marginal-boundary,obj:thm:uniform-mixed-upper,obj:prop:binary-alphabet-parametric-rate,obj:lem:chebyshev-factorial-certificate,obj:lem:fixed-sample-hybrid-transfer,obj:lem:finite-alternation-duality,obj:lem:parametric-label-floor,obj:lem:poisson-inverse-marginal-energy-scalar,obj:lem:poisson-prefix-count-law,obj:lem:inverse-count-poisson-prefix-transfer,obj:lem:shifted-chebyshev-quotient-coefficient-bound,obj:lem:annotation-frontier-strict-improvement,obj:lem:common-marginal-recipe-target-concentration,obj:lem:common-marginal-recipe-scale-tail,obj:lem:common-marginal-canonical-certificate-family,obj:lem:calibrated-envelope-absorption,obj:lem:finite-side-measurable-comparison,obj:lem:hybrid-light-branch-mean-bound,obj:lem:poisson-inverse-arm-mean-formula,obj:lem:heavy-light-second-moment-sum,obj:lem:hybrid-heavy-branch-mean-bound,obj:lem:calibrated-hybrid-variance-bound,obj:lem:calibrated-hybrid-mean-bias,obj:lem:finite-product-padding-identity,obj:lem:finite-side-cover-approximate-decision,obj:thm:common-marginal-converse,obj:thm:sharp-annotation-frontier,obj:lem:poisson-inverse-count-risk,obj:thm:inverse-count-baseline,obj:lem:poisson-hybrid-risk,obj:lem:explicit-chebyshev-calibration,obj:lem:common-marginal-uniform-intensity,obj:lem:finite-side-information-convergence,obj:lem:common-marginal-recipe-transfer,obj:lem:fixed-sample-common-marginal-transfer,obj:thm:known-marginal-limit}.
\item \textbf{(Integer.)} Let \(n\ge1\) be an integer.
\end{itemize}
Define
\[
b_n:=1+\lceil\log_2 n\rceil .
\]
Equivalently, in the notation of \cref{obj:synth_7}, \(b_n=h_n\).
\end{definitionv}

\begin{definitionv}{synth_3}[Clipping to the unit interval]
For \(z\in\mathbb R\), define
\[
\operatorname{clip}_{[-1,1]}(z)
=
\max\{-1,\min\{1,z\}\}.
\]
Equivalently, \(\operatorname{clip}_{[-1,1]}(z)\) is the Euclidean projection of \(z\) onto the interval \([-1,1]\).
\end{definitionv}

\begin{algorithmv}{def:mixed-estimator-handle}[Mixed estimator definition]
For \(n,m,d\in\mathbb N\) and a real overlap constant \(\epsilon\), define
\[
\widehat\tau^{\mathrm{mix}}_{n,m,d}: \mathcal T_{n,m,d}
\]
as the following map from
\[
\bigl([d]\times\{0,1\}\times\{0,1\}\bigr)^n
\times
\bigl([d]\times\{0,1\}\bigr)^m
\]
to \(\mathbb R\). Put
\[
M_0=\lfloor n/3\rfloor,\qquad
n_p=\left\lfloor\frac{n-M_0}{2}\right\rfloor,\qquad
n_f=n-M_0-n_p,
\]
\[
m_p=\lfloor m/2\rfloor,\qquad
m_f=m-m_p,\qquad
M_p=n_p+m_p,\qquad
M_f=n_f+m_f,
\]
and
\[
u=\frac{M_0}{8},\qquad
t_p=\frac{M_p}{8},\qquad
t=\frac{M_f}{8}.
\]
Let \(b_n\) denote the binary-length calibration quantity used here, and set
\[
L=L(n),\qquad B=B(n,m,\epsilon),\qquad k_0=k_0(n,m,\epsilon),
\]
with \(\bar\epsilon=\bar\epsilon(\epsilon)\), \(A_\star=7056\), and \(c^\circ=1/256\). The calibration predicate \(\mathsf C_\epsilon(n,m)\) holds exactly when all of the following finite inequalities hold:
\begin{itemize}
\item \textbf{(Sample size.)} \(24\le n\).
\item \textbf{(Block proportions.)}
\[
\frac1{32}\le\frac un\le\frac1{24},\qquad
\frac1{32}\le\frac{t_p}{n+m}\le\frac3{32},\qquad
\frac1{32}\le\frac t{n+m}\le\frac3{32},\qquad
\frac13\le\frac t{t_p}\le3.
\]
\item \textbf{(Degree and scale.)}
\[
\frac{c^\circ b_n}{2}\le L,\qquad
8\le t_pB,\qquad
L\le tB.
\]
\item \textbf{(Polynomial size.)}
\[
L^4A_\star^{2L}\le n,\qquad
L^6A_\star^{2L}\le n.
\]
\item \textbf{(Tail calibration.)}
\[
15L+12b_n\le\frac{t_pB}{4},\qquad
n^{12}\le\left(1+\bar\epsilon\frac t{t_p}\right)^{k_0},\qquad
12b_n\le\bar\epsilon tB.
\]
\end{itemize}
If \(\mathsf C_\epsilon(n,m)\) fails, set
\[
\widehat\tau^{\mathrm{mix}}_{n,m,d}(\mathcal L_n,\mathcal U_m)=0.
\]
If \(\mathsf C_\epsilon(n,m)\) holds, set
\[
\widehat\tau^{\mathrm{mix}}_{n,m,d}(\mathcal L_n,\mathcal U_m)
=
\sum_{r_0=0}^{M_0}
\sum_{r_p=0}^{M_p}
\sum_{r_f=0}^{M_f}
\Pr\{\operatorname{Poi}(u)=r_0\}
\Pr\{\operatorname{Poi}(t_p)=r_p\}
\Pr\{\operatorname{Poi}(t)=r_f\}
\,\Phi_\epsilon(\mathcal L_n,\mathcal U_m;r_0,r_p,r_f),
\]
where
\[
\Phi_\epsilon(\mathcal L_n,\mathcal U_m;r_0,r_p,r_f)
=
\operatorname{clip}_{[-1,1]}\left(\sum_{x\in[d]} Z_x\right).
\]
For each \(x\in[d]\), let \(J_x\) be the pilot-prefix count in cell \(x\), let \(S_{ax}\) be the outcome-block prefix count with covariate \(x\), treatment \(a\), and outcome \(1\), and let \(K_{ax}\) be the factorial-pool prefix count with covariate \(x\) and treatment \(a\). Then
\[
Z_x
=
\begin{cases}
\dfrac{S_{1x}}u\,\widehat W_{1,L,x}
-
\dfrac{S_{0x}}u\,\widehat W_{0,L,x},
& J_x\le k_0,\\[1.2em]
\dfrac{S_{1x}}u\,\dfrac{K_{0x}+K_{1x}+1}{K_{1x}+1}
-
\dfrac{S_{0x}}u\,\dfrac{K_{0x}+K_{1x}+1}{K_{0x}+1},
& k_0<J_x.
\end{cases}
\]
Here \(\widehat W_{a,L,x}\) is the factorial lift, at degree \(L\), scale \(B\), and intensity \(t\), of the light-cell polynomial weight
\[
W_{a,L}(s_0,s_1)=\frac{s_0+s_1}{B}G_L(s_a/B),
\]
where
\[
H_L(z)=\frac{1-T_L(1-2z)}{2L^2},\qquad
E_L(z)=\frac{H_L(z)}z,\qquad
G_L(z)=\frac{1-E_L(z)}z,
\]
with polynomial continuation at zero.
\end{algorithmv}

\begin{theoremv}{thm:uniform-mixed-upper}[Uniform mixed upper bound]
Fix an overlap constant \(\epsilon\) with \(0<\epsilon<1/2\).  There exists a constant \(C_\epsilon>0\) such that, for every \(n\ge 1\), \(m\ge 0\), and \(d\ge 2\), the total computable hybrid estimator \(\widehat\tau^{\mathrm{mix}}_{n,m,d}\) of \cref{obj:def:mixed-estimator-handle} satisfies
\[
\sup_{P\in\mathcal M_{d,\epsilon}}
E_{P^{\otimes n}\otimes P_{XA}^{\otimes m}}
\!\left[
\left\{\widehat\tau^{\mathrm{mix}}_{n,m,d}-\tau(P)\right\}^2
\right]
\le
C_\epsilon\, r_\epsilon(n,m,d),
\]
where \(\tau(P)\) is the average treatment effect functional in \cref{obj:def:ate-functional}, and \(r_\epsilon(n,m,d)\) is the annotation-frontier rate in \cref{obj:def:frontier-rate-handle}.
\end{theoremv}

\begin{definitionv}{synth_22}[Common-marginal finite-atom construction]
Fix \(0<\epsilon<1/2\), an integer \(L\ge2\), scales \(B>0\) and \(a>0\), and write
\[
\kappa=\frac{1-2\epsilon}{\epsilon},\qquad I=[a/\kappa,B].
\]
Let \(\sigma\) be a finite signed measure on \(I\) with total variation one, let
\(h=d\sigma/d|\sigma|\), and define a probability law \(\nu\) on \(\{0\}\cup I\) by
\[
\nu(dp)=\frac{a}{p+a}\,|\sigma|(dp)\quad(p\in I),
\qquad
\nu(\{0\})=1-\int_I\frac{a}{p+a}\,|\sigma|(dp).
\]
For an integer \(k\le d-1\), the finite-atom construction has atom-generator law
\[
\Gamma=\nu^{\otimes k}
\]
on \(w=(p_1,\ldots,p_k)\).  Given \(w\), form \(k\) rare cells and one reservoir cell.  For each rare cell with \(p_i>0\), set
\[
s_i=\epsilon(p_i+a),\qquad e_i=\frac{s_i}{p_i},
\]
put \(\mu_{0i}=0\), and set the treated outcome regression in branch \(b\in\{0,1\}\) to
\[
\mu_{1i}^{(b)}=\frac{1+(2b-1)h(p_i)}{2}.
\]
At atoms with \(p_i=0\), assign the same legal propensity and outcome values in both branches.  The reservoir has raw mass
\[
p_*(w)=1-kE_\nu[p],
\]
propensity \(1/2\), and outcome regressions equal to zero.  Alphabet cells not used by the \(k\) rare cells or the reservoir have mass zero.

The raw total mass is
\[
S(w)=p_*(w)+\sum_{i=1}^k p_i.
\]
The branch law \(P_b(w)\) is obtained by dividing every raw cell mass by \(S(w)\) and keeping the propensities and outcome regressions just specified.  The corresponding raw target is
\[
r_b(w)=\sum_{i=1}^k p_i\,\mu_{1i}^{(b)},
\]
so the normalized target value of the generated law is \(r_b(w)/S(w)\).  The branch priors are the pushforwards
\[
\Pi_b=\Gamma\circ P_b^{-1},\qquad b\in\{0,1\}.
\]
Thus the two branches use the same generated rare-cell masses, reservoir mass, propensities, normalization factor, and treatment--covariate table, while differing only in the branch sign of the treated rare-cell outcome regressions.
\end{definitionv}

\begin{definitionv}{def:common-marginal-prior-handle}[Common-marginal prior handle]
Fix \(n,m,d\in\mathbb N\), \(\epsilon\in\mathbb R\), and a calibration consisting of constants
\[
C_\epsilon,\quad C_\epsilon^\circ,\quad b_0,\quad b_\epsilon,\quad
\gamma_\epsilon,\quad c_\epsilon^\circ,\quad c_\epsilon,\quad C_d
\]
with all eight constants positive and \(c_\epsilon\le 1/4\).  A common-marginal prior handle at \((n,m,d)\) consists of integers \(L,k\), real numbers \(B,a\), and a calibrated recipe supplying two induced priors, denoted \(\Pi_0(H)\) and \(\Pi_1(H)\), on laws on \([d]\times\{0,1\}^2\), subject to
\[
L=\left\lceil C_\epsilon^\circ \log(en)\right\rceil,\qquad
B=\frac{b_0L}{n+m},\qquad
a=\frac{\gamma_\epsilon B}{L^2},
\]
and
\[
k=\min\left\{d-1,\left\lfloor c_\epsilon^\circ(n+m)L\right\rfloor\right\}.
\]
The recipe uses the finite-atom construction exactly when
\[
\min\left\{1,\frac{d^2}{(n+m)^2\log(en)^2}\right\}
>
\frac{C_\epsilon}{n}.
\]
For a handle \(H\) and a branch \(b\in\{0,1\}\), the common-marginal prior is the branch-indexed measure
\[
\Pi_b(H)=
\begin{cases}
\Pi_1(H), & b=1,\\
\Pi_0(H), & b=0.
\end{cases}
\]
\end{definitionv}

\begin{theoremv}{thm:common-marginal-converse}[Common marginal converse]
Let \(\epsilon\), the overlap parameter, satisfy \(0<\epsilon<1/2\).  Then there exist a constant \(c_\epsilon>0\) and a collection of \(\epsilon\)-dependent calibration constants as in \cref{obj:def:common-marginal-prior-handle} such that, for every \(n,m,d\in\mathbb N\) with \(n\ge1\), \(m\ge0\), and \(d\ge2\), the regime-dependent handle in \cref{obj:def:common-marginal-prior-handle} supplies priors \(\Pi_0\) and \(\Pi_1\) satisfying
\[
\mathcal L_{\Pi_0}(P_{XA})=\mathcal L_{\Pi_1}(P_{XA})
\]
and
\[
c_\epsilon\, r_\epsilon(n,m,d)\le R_\epsilon(n,m,d),
\]
where \(r_\epsilon(n,m,d)\) is the frontier rate in \cref{obj:def:frontier-rate-handle} and \(R_\epsilon(n,m,d)\) is the minimax risk in \cref{obj:def:minimax-risk}.
\end{theoremv}

\begin{theoremv}{thm:sharp-annotation-frontier}[Sharp annotation frontier]
For every overlap parameter \(\epsilon\) satisfying \(0<\epsilon<1/2\), there exist constants \(c_\epsilon\) and \(C_\epsilon\), depending only on \(\epsilon\), with
\[
0<c_\epsilon\le C_\epsilon,
\]
such that the following assertions hold.

\begin{itemize}
\item \textbf{(Finite-sample bracket.)} For every \(n\ge1\), \(m\ge0\), and \(d\ge2\),
\[
c_\epsilon r_\epsilon(n,m,d)
\le R_\epsilon(n,m,d)
\le C_\epsilon r_\epsilon(n,m,d),
\]
where, with \(N=n+m\) and \(\ell_n=\log(en)\),
\[
r_\epsilon(n,m,d)
=
\min\!\left\{1,\frac1n+\frac{d^2}{N^2\ell_n^2}\right\}.
\]

\item \textbf{(Unannotated endpoint.)} For every \(n\ge1\) and every \(d\ge0\),
\[
r_\epsilon(n,0,d)
=
\min\!\left\{1,\frac1n+\frac{d^2}{n^2\log^2(en)}\right\}.
\]

\item \textbf{(Known-marginal limit.)} For every fixed \(n\ge1\) and \(d\ge2\),
\[
\lim_{m\to\infty} R_\epsilon(n,m,d)
=
R_\epsilon^{\mathrm{known}}(n,d).
\]

\item \textbf{(Sequence regimes.)} For any sequences \((m_n)_{n\ge1}\) and \((d_n)_{n\ge1}\) of nonnegative integers satisfying \(d_n\ge2\) for every \(n\ge1\), set \(N_n=n+m_n\) and \(\ell_n=\log(en)\). Then
\[
R_\epsilon(n,m_n,d_n)\to0
\quad\Longleftrightarrow\quad
\frac{d_n}{N_n\ell_n}\to0,
\]
\[
R_\epsilon(n,m_n,d_n)=O(n^{-1})
\quad\Longleftrightarrow\quad
d_n=O\!\left(\frac{N_n\ell_n}{\sqrt n}\right),
\]
and
\[
R_\epsilon(n,m_n,d_n)
=
o\!\left(R_\epsilon(n,0,d_n)\right)
\]
if and only if both
\[
\frac{d_n}{\sqrt n\,\ell_n}\to\infty
\]
and
\[
\frac{d_n/(N_n\ell_n)}
{\min\{1,d_n/(n\ell_n)\}}
\to0.
\]
\end{itemize}
\end{theoremv}

\begin{theoremv}{thm:known-marginal-limit}[Known marginal limit]
Fix an overlap parameter \(\epsilon\), a labeled sample size \(n\), and a finite covariate alphabet size \(d\). Suppose
\begin{itemize}
\item \textbf{(Overlap.)} \(0<\epsilon<1/2\).
\item \textbf{(Labeled size.)} \(n\ge 1\).
\item \textbf{(Alphabet size.)} \(d\ge 2\).
\end{itemize}
Then the minimax mean-squared risk \(R_\epsilon(n,m,d)\) from \cref{obj:def:minimax-risk} converges, as the auxiliary sample size \(m\to\infty\), to the known-marginal minimax risk \(R_\epsilon^{\mathrm{known}}(n,d)\) from \cref{obj:def:known-marginal-risk}:
\[
\lim_{m\to\infty} R_\epsilon(n,m,d)
=
R_\epsilon^{\mathrm{known}}(n,d).
\]
\end{theoremv}

\begin{definitionv}{synth_5}[Poisson-prefix inverse-count estimator]
For \(n\ge1\), \(m\ge0\), and \(d\ge2\), put
\[
M_0=\lfloor n/3\rfloor,\qquad
n_p=\left\lfloor\frac{n-M_0}{2}\right\rfloor,\qquad
n_f=n-M_0-n_p,\qquad
m_p=\lfloor m/2\rfloor,\qquad
m_f=m-m_p,
\]
and \(u=M_0/8\). Given a labeled sample \(\mathcal L_n\) and an auxiliary sample \(\mathcal U_m\), define \(\widehat\tau^{\mathrm{IC,prefix}}_{n,m,d}(r_0,r_f)=0\) when \(r_0>M_0\) or \(r_f>n_f+m_f\). For \(0\le r_0\le M_0\) and \(0\le r_f\le n_f+m_f\), let \(S_{ax}(r_0)\) be the number of records among the first \(r_0\) labeled observations with \((X,A,Y)=(x,a,1)\). Let \(K_{ax}(r_f)\) be the number of records with \((X,A)=(x,a)\) in the factorial prefix formed from the last \(n_f\) labeled records and the last \(m_f\) auxiliary records, taking the first \(r_f\) records of that pooled factorial block in its fixed order. Define
\[
\widehat\tau^{\mathrm{IC,prefix}}_{n,m,d}(r_0,r_f)
=
\operatorname{clip}_{[-1,1]}
\left[
\sum_{x=1}^d
\left\{
\frac{S_{1x}(r_0)}{u}
\frac{K_{0x}(r_f)+K_{1x}(r_f)+1}{K_{1x}(r_f)+1}
-
\frac{S_{0x}(r_0)}{u}
\frac{K_{0x}(r_f)+K_{1x}(r_f)+1}{K_{0x}(r_f)+1}
\right\}
\right].
\]
\end{definitionv}

\begin{theoremv}{thm:inverse-count-baseline}[Inverse-count baseline bound]
Fix an overlap constant \(\epsilon\) with \(0<\epsilon<1/2\).  Then there is a constant \(C_\epsilon>0\) such that, for every \(n\ge1\), \(m\ge0\), and \(d\ge2\), the clipped derandomized inverse-count estimator defined for each sample \(Z=(\mathcal L_n,\mathcal U_m)\) by
\[
\widehat\tau^{\mathrm{IC}}_{n,m,d}(Z)
=
\sum_{r_0=0}^{M_0}
\sum_{r_f=0}^{n_f+m_f}
\mathbb P\{\operatorname{Poi}(M_0/8)=r_0\}
\mathbb P\{\operatorname{Poi}((n_f+m_f)/8)=r_f\}
\,\widehat\tau^{\mathrm{IC,prefix}}_{n,m,d}(Z;r_0,r_f)
\]
is measurable.  Here \(M_0=\lfloor n/3\rfloor\), \(n_p=\lfloor(n-M_0)/2\rfloor\), \(n_f=n-M_0-n_p\), \(m_p=\lfloor m/2\rfloor\), and \(m_f=m-m_p\) are the block sizes from \cref{obj:def:mixed-estimator-handle}, and \(\widehat\tau^{\mathrm{IC,prefix}}_{n,m,d}(Z;r_0,r_f)\) is the sample-and-prefix-count statistic of \cref{obj:synth_5}.

Moreover,
\[
\sup_{P\in\mathcal M_{d,\epsilon}}
\mathbb E_{P^{\otimes n}\otimes P_{XA}^{\otimes m}}
\left[
\left\{
\widehat\tau^{\mathrm{IC}}_{n,m,d}-\tau(P)
\right\}^2
\right]
\le
C_\epsilon
\min\left\{
1,\frac1n+\frac{d^2}{(n+m)^2}
\right\}.
\]
Here \(\mathcal M_{d,\epsilon}\) is the overlap class in \cref{obj:def:model-class}, \(\tau(P)\) is the average treatment effect functional in \cref{obj:def:ate-functional}, and the expectation is under the annotation experiment of \cref{obj:def:annotation-experiment}.  Consequently, for the minimax risk in \cref{obj:def:minimax-risk},
\[
R_\epsilon(n,m,d)
\le
C_\epsilon
\min\left\{
1,\frac1n+\frac{d^2}{(n+m)^2}
\right\}.
\]
\end{theoremv}

\begin{theoremv}{prop:binary-alphabet-parametric-rate}[Binary parametric rate]
Let \(0<\epsilon<1/2\). For the minimax risk \(R_\epsilon(n,m,d)\) defined in \cref{obj:def:minimax-risk}, at binary alphabet size \(d=2\) there exist constants \(c_\epsilon,C_\epsilon\in\mathbb R\) with \(0<c_\epsilon\le C_\epsilon\) such that, for every integer \(n\ge1\) and every integer \(m\ge0\),
\[
\frac{c_\epsilon}{n}\le R_\epsilon(n,m,2)\le \frac{C_\epsilon}{n}.
\]
\end{theoremv}

\begin{lemmav}{lem:poisson-inverse-marginal-energy-scalar}[Marginal energy insertion]
Let \(u,p,w,\epsilon,\lambda_0,\lambda_1,q_0,q_1,r_{2,0},r_{2,1},r_{4,0},r_{4,1}\in\mathbb R\).  Suppose that
\begin{itemize}
\item \textbf{(Basic scales.)} \(u>0\), \(p\ge0\), \(0<\epsilon\le1\), and \(w\ge0\).
\item \textbf{(Outcome masses.)} \(0\le q_0\le p\) and \(0\le q_1\le p\).
\item \textbf{(Marginal intensities.)} For each \(a\in\{0,1\}\), \(\epsilon w\le\lambda_a\le w\).
\item \textbf{(Second-moment rates.)} \(r_{2,0},r_{2,1}\ge0\), \((1+\lambda_0)^2r_{2,0}\le8\), and \((1+\lambda_1)^2r_{2,1}\le8\).
\item \textbf{(Fourth-moment rates.)} \(r_{4,0},r_{4,1}\ge0\), \((1+\lambda_0)^4r_{4,0}\le96\), and \((1+\lambda_1)^4r_{4,1}\le96\).
\end{itemize}
Put
\[
c_{\mathrm{mar}}=8\epsilon^{-2}+96(\epsilon^{-2}+\epsilon^{-3}).
\]
Then
\[
\lambda_1\bigl\{(q_1u^{-1}+q_1^2)(\lambda_0+\lambda_0^2)r_{4,1}
 +(q_0u^{-1}+q_0^2)r_{2,0}\bigr\}
+\lambda_0\bigl\{(q_0u^{-1}+q_0^2)(\lambda_1+\lambda_1^2)r_{4,0}
 +(q_1u^{-1}+q_1^2)r_{2,1}\bigr\}
\le
(pu^{-1}+p^2)\frac{2c_{\mathrm{mar}}}{1+w}.
\]
\end{lemmav}

\begin{definitionv}{synth_2}[Cell arm mass]
For a law \(P\) on \([d]\times\{0,1\}^2\), a covariate cell \(x\in[d]\), and a treatment arm \(a\in\{0,1\}\), define the treatment-covariate cell mass
\[
s_{ax}
=
\sum_{y\in\{0,1\}} P(X=x,A=a,Y=y).
\]
\end{definitionv}

\begin{definitionv}{synth_1}[Outcome-marked cell mass]
For an observed-data law \(P\) on the finite record space \([d]\times\{0,1\}\times\{0,1\}\), define the outcome-marked cell mass \(q_{ax}\) for \(x\in[d]\) and \(a\in\{0,1\}\) by
\[
q_{ax}
=
P(X=x,A=a,Y=1).
\]
Equivalently, \(q_{ax}\) is the probability mass assigned by \(P\) to the complete atom \((x,a,1)\).
\end{definitionv}

\begin{lemmav}{lem:poisson-inverse-count-risk}[Poisson inverse-count risk]
Fix an overlap constant \(\epsilon\) with \(0<\epsilon<1/2\).  There is a constant \(C=C_\epsilon>0\) such that, for every \(d\), every \(u,t>0\), and every observed-data law \(P\in\mathcal M_{d,\epsilon}\) as in \cref{obj:def:model-class}, the following holds.

In the independent Poisson count experiment with independent counts
\[
S_{ax}\sim \operatorname{Poi}(u q_{ax}),
\qquad
K_{ax}\sim \operatorname{Poi}(t s_{ax}),
\qquad a\in\{0,1\},\ x\in[d],
\]
define
\[
Z^H
=
\sum_{x=1}^d
\left\{
\frac{S_{1x}}{u}\frac{K_{0x}+K_{1x}+1}{K_{1x}+1}
-
\frac{S_{0x}}{u}\frac{K_{0x}+K_{1x}+1}{K_{0x}+1}
\right\}.
\]
Writing \(\tau(P)\) for the average treatment effect functional in \cref{obj:def:ate-functional},
\[
\left|E_P Z^H-\tau(P)\right|
\le
C\min\{1,d/t\},
\qquad
\operatorname{Var}_P(Z^H)
\le
C\left(u^{-1}+t^{-1}\right).
\]
\end{lemmav}

\begin{lemmav}{lem:poisson-prefix-count-law}[Poisson prefix count law]
Let \(P\) be a probability law on \([d]\times\{0,1\}^2\), and let \(u,t\ge 0\).  Write
\[
q_{ax}=P(x,a,1),
\qquad
s_{ax}=P(x,a,0)+P(x,a,1).
\]
Under the prefix law that independently draws an outcome-marked prefix with \(\operatorname{Poi}(u)\) records from \(P\) and a treatment--covariate prefix with \(\operatorname{Poi}(t)\) records from the marginal \(P_{XA}\), define
\[
S_{ax}=\#\{\text{outcome-marked prefix records equal to }(x,a,1)\},
\qquad
K_{ax}=\#\{\text{treatment--covariate prefix records equal to }(x,a)\}.
\]
Then the pushforward of the prefix law under the map sending the two prefixes to the array \((S_{ax},K_{ax})_{x\in[d],\,a\in\{0,1\}}\) is the product measure
\[
\bigotimes_{x\in[d]}\bigotimes_{a\in\{0,1\}}
\left\{\operatorname{Poi}(u q_{ax})\otimes \operatorname{Poi}(t s_{ax})\right\}.
\]
Equivalently, the coordinates are mutually independent, with
\[
S_{ax}\sim \operatorname{Poi}(u q_{ax}),
\qquad
K_{ax}\sim \operatorname{Poi}(t s_{ax})
\]
for every \(x\in[d]\) and \(a\in\{0,1\}\).
\end{lemmav}

\begin{definitionv}{synth_19}[Poisson-prefix law]
For \(P\in\mathcal M_{d,\epsilon}\) and intensities \(u,t\ge0\), let \(\mathbb P^{\mathrm{pref}}_{P,u,t}\) denote the following probability law. Draw two prefix lengths independently,
\[
R_0\sim \operatorname{Poi}(u),\qquad R_f\sim \operatorname{Poi}(t).
\]
Conditional on \((R_0,R_f)=(r_0,r_f)\), draw an outcome prefix
\[
((X_i,A_i,Y_i))_{i=1}^{r_0}\sim P^{\otimes r_0}
\]
and, independently, a factorial treatment--covariate prefix
\[
((\widetilde X_j,\widetilde A_j))_{j=1}^{r_f}\sim P_{XA}^{\otimes r_f}.
\]
Thus the two prefixes are independent, the outcome prefix has complete-record sampling law \(P\), and the factorial prefix has treatment--covariate sampling law \(P_{XA}\). Under \(\mathbb P^{\mathrm{pref}}_{P,u,t}\), the associated counts
\[
S_{ax}=\sum_{i=1}^{R_0}\mathbf 1\{X_i=x,A_i=a,Y_i=1\},
\qquad
K_{ax}=\sum_{j=1}^{R_f}\mathbf 1\{\widetilde X_j=x,\widetilde A_j=a\}
\]
are independent across the two prefixes, with
\[
S_{ax}\sim\operatorname{Poi}(u q_{ax}),\qquad
K_{ax}\sim\operatorname{Poi}(t s_{ax}).
\]
\end{definitionv}

\begin{lemmav}{lem:inverse-count-poisson-prefix-transfer}[Poisson prefix transfer]
Let \(n,m,d\) be nonnegative integers, let \(\epsilon\in\mathbb R\), and let \(P\in\mathcal M_{d,\epsilon}\) as in \cref{obj:def:model-class}.  Define
\[
M_0=\bigl\lfloor n/3\bigr\rfloor,\qquad
n_p=\bigl\lfloor (n-M_0)/2\bigr\rfloor,\qquad
n_f=n-M_0-n_p,\qquad
m_f=m-\bigl\lfloor m/2\bigr\rfloor,
\]
and set \(u=M_0/8\) and \(t=(n_f+m_f)/8\).  Let \(\widehat\tau^{\mathrm{IC}}_{n,m,d}\) be the clipped derandomized inverse-count estimator obtained by averaging, over \(0\le r_0\le M_0\) and \(0\le r_f\le n_f+m_f\), the finite-prefix inverse-count output with weights
\[
\Pr\{\operatorname{Poi}(u)=r_0\}\Pr\{\operatorname{Poi}(t)=r_f\}.
\]
Let \(\widehat\tau^{\mathrm{IC,pref}}_{u,t}\) be the clipped inverse-count statistic computed from two independent Poisson prefixes: an outcome-marked prefix with law induced by \(P\) and length \(\operatorname{Poi}(u)\), and a treatment--covariate prefix with law \(P_{XA}\) and length \(\operatorname{Poi}(t)\).  Then, for the annotation experiment in \cref{obj:def:annotation-experiment} and the ATE functional \(\tau(P)\) in \cref{obj:def:ate-functional},
\[
E_{P^{\otimes n}\otimes P_{XA}^{\otimes m}}
\!\left[
  \bigl(\widehat\tau^{\mathrm{IC}}_{n,m,d}-\tau(P)\bigr)^2
\right]
\le
E_{\mathbb P^{\mathrm{pref}}_{P,u,t}}
\!\left[
  \bigl(\widehat\tau^{\mathrm{IC,pref}}_{u,t}-\tau(P)\bigr)^2
\right]
+4\Bigl[
  \Pr\{\operatorname{Poi}(u)>M_0\}
  +\Pr\{\operatorname{Poi}(t)>n_f+m_f\}
\Bigr].
\]
\end{lemmav}

\begin{definitionv}{synth_4}[Factorial lift of the light-cell polynomial]
Fix \(a\in\{0,1\}\), an integer \(L\ge2\), and constants \(B,t>0\). Let
\[
W_{a,L}(s_0,s_1)
=
\frac{s_0+s_1}{B}G_L(s_a/B)
=
\sum_{r_0,r_1\ge0}c_{a,r_0,r_1}s_0^{r_0}s_1^{r_1},
\]
where \(G_L\) is the Chebyshev correction polynomial and only finitely many coefficients \(c_{a,r_0,r_1}\) are nonzero. For integer counts \(K_0,K_1\ge0\), define the factorial lift
\[
\widehat W_{a,L}(K_0,K_1)
=
\sum_{r_0,r_1\ge0}
c_{a,r_0,r_1}
\frac{(K_0)_{r_0}(K_1)_{r_1}}{t^{r_0+r_1}},
\]
with \((k)_r=k(k-1)\cdots(k-r+1)\) and \((k)_0=1\). Thus each monomial of \(W_{a,L}\) is replaced by the corresponding unbiased Poisson factorial-moment statistic.
\end{definitionv}

\begin{lemmav}{lem:chebyshev-factorial-certificate}[Chebyshev factorial certificate]
There is a universal constant \(A>1\) such that, for every integer \(L\ge2\), the polynomial continuations
\[
E_L(z)=H_L(z)/z,
\qquad
G_L(z)=\{1-E_L(z)\}/z
\]
satisfy the following properties:
\begin{itemize}
\item \textbf{(Degree.)} The polynomial \(G_L\) has degree at most \(L-2\).
\item \textbf{(Chebyshev certificate.)} For every \(z\in[0,1]\),
\[
0\le E_L(z)\le 1,
\qquad
zG_L(z)+E_L(z)=1,
\]
and, whenever \(z>0\),
\[
E_L(z)\le (L^2z)^{-1}.
\]
\item \textbf{(Factorial-lift moment.)} For every \(a\in\{0,1\}\) and every \(t,B,s_0,s_1\in\mathbb R\) with
\[
t>0,\qquad B>0,\qquad s_0\ge0,\qquad s_1\ge0,\qquad L\le tB,
\]
if \(K_0\sim\operatorname{Poi}(ts_0)\) and \(K_1\sim\operatorname{Poi}(ts_1)\) are independent, then the factorial lift \(\widehat W_{a,L}\) satisfies
\[
\mathbb E\!\left[\widehat W_{a,L}(K_0,K_1)^2\right]
\le
A^L\left(1+\frac{s_0+s_1}{B}\right)^{2L}.
\]
Moreover, if \(s_0+s_1\le B\), then
\[
\mathbb E\!\left[\widehat W_{a,L}(K_0,K_1)^2\right]\le A^L.
\]
\end{itemize}
\end{lemmav}

\begin{lemmav}{lem:calibrated-envelope-absorption}[Calibrated envelope absorption]
Let \(\epsilon>0\) and \(C_0>0\). There is a constant \(C>0\), depending on \(\epsilon\) and \(C_0\), such that the following holds.

\begin{itemize}
\item \textbf{(Sample sizes.)} \(n,m,d\) are integers with \(n\ge 1\), \(m\ge 0\), and \(d\ge 2\).
\item \textbf{(Calibration.)} The finite calibration predicate \(\mathsf C_\epsilon(n,m)\) from \cref{obj:def:mixed-estimator-handle} holds.
\item \textbf{(Estimator scales.)} With the block sizes and scales from \cref{obj:def:mixed-estimator-handle},
\[
M_0=\lfloor n/3\rfloor,\qquad
n_p=\lfloor (n-M_0)/2\rfloor,\qquad
n_f=n-M_0-n_p,\qquad
m_f=m-\lfloor m/2\rfloor,
\]
and
\[
u=M_0/8,\qquad
t=(n_f+m_f)/8,\qquad
L=L(n),\qquad
B=B(n,m,\epsilon),\qquad
A_\star=7056.
\]
Let \(N=n+m\) and \(\ell_n=\log(en)\), as in \cref{obj:def:frontier-rate-handle}.
\end{itemize}

Then
\[
4A_\star^L
\left\{
\frac{\min\{1,dB\}}{u}+dB^2
\right\}
+C_0\left(u^{-1}+t^{-1}\right)
+\frac{8dB^2}{\epsilon^2L^4}
+\frac{4}{\epsilon t}
+414\,n^{-12}
+\left(
\frac{2dB}{\epsilon L^2}
+6n^{-12}
\right)^2
\le
C\left(
\frac{1}{n}
+\frac{d^2}{N^2\ell_n^2}
\right).
\]
\end{lemmav}

\begin{lemmav}{lem:hybrid-light-branch-mean-bound}[Light branch mean bound]
Let \(P\in\mathcal M_{d,\epsilon}\) be a law in the finite model class of \cref{obj:def:model-class}.  Let \(u,t,B>0\) and let \(L\ge 2\) be an integer.  For each \(x\in[d]\), write
\[
p_x=P(X=x),\qquad
q_{ax}=P(X=x,A=a,Y=1),\qquad
s_{ax}=P(X=x,A=a).
\]
For each \(x\in[d]\), let \(S_{0x},S_{1x},K_{0x},K_{1x}\) be independent Poisson counts with
\[
S_{ax}\sim \operatorname{Poi}(u q_{ax}),
\qquad
K_{ax}\sim \operatorname{Poi}(t s_{ax}),
\qquad a\in\{0,1\}.
\]
Using the light factorial-lift weights and heavy inverse-count weights from \cref{obj:def:mixed-estimator-handle}, set
\[
P_x
=
\frac{S_{1x}}{u}\widehat W_{1,L,x}
-
\frac{S_{0x}}{u}\widehat W_{0,L,x},
\qquad
H_x
=
\frac{S_{1x}}{u}\widehat W^H_{1,x}
-
\frac{S_{0x}}{u}\widehat W^H_{0,x},
\]
where the light weights use degree \(L\), scale \(B\), and intensity \(t\).  Then
\[
\sum_{x:\,p_x\le B}
\left(\mathbb E P_x-\mathbb E H_x\right)^2
\le
\frac{8dB^2}{\epsilon^2 L^4}
+
\frac{4}{\epsilon t}.
\]
\end{lemmav}

\begin{lemmav}{lem:poisson-inverse-arm-mean-formula}[Inverse-count arm mean]
Let \(P\) be an observed-data law on \([d]\times\{0,1\}^2\), let \(x\in[d]\), and let \(a\in\{0,1\}\). Write \(p_x\) for the covariate-cell mass, \(s_{ax}\) for the treatment-covariate cell mass, \(q_{ax}\) for the outcome-marked mass, and \(\mu_{ax}:=q_{ax}/s_{ax}\) for the totalized outcome regression. If \(u>0\) and \(t>0\), and if \(K_{ax}\sim\operatorname{Poi}(t s_{ax})\) and \(K_{1-a,x}\sim\operatorname{Poi}(t s_{1-a,x})\) are independent, then
\[
\frac{q_{ax}}{u/u}\,
\mathbb E\!\left[1+\frac{K_{1-a,x}}{K_{ax}+1}\right]
=
p_x\mu_{ax}-(p_x-s_{ax})\mu_{ax}\exp(-t s_{ax}).
\]
\end{lemmav}

\begin{lemmav}{lem:heavy-light-second-moment-sum}[Heavy light second moment]
Let \(d\) be a finite alphabet size, let \(P\) be an observed-data law on \([d]\times\{0,1\}^2\), and let \(n,m\) be integers. Suppose that:
\begin{itemize}
\item \textbf{(Positive overlap parameter.)} The overlap parameter satisfies \(\epsilon>0\).
\item \textbf{(Sample size.)} The labeled sample size satisfies \(n\ge1\).
\item \textbf{(Calibration.)} The finite calibration predicate \(\mathsf C_\epsilon(n,m)\) from \cref{obj:def:mixed-estimator-handle} holds.
\end{itemize}
Set
\[
u=\frac{M_0}{8},\qquad
t_p=\frac{M_p}{8},\qquad
t=\frac{M_f}{8},\qquad
L=L(n),\qquad
B=B(n,m,\epsilon),\qquad
k_0=k_0(n,m,\epsilon),
\]
with these quantities defined as in \cref{obj:def:mixed-estimator-handle}. For each \(x\in[d]\), let
\[
\pi_x
=
\mathbb P\{J_x\le k_0\},
\]
where, conditionally on \(P\), the pilot counts have independent components
\(J_{ax}\sim\operatorname{Poi}(t_p s_{ax})\), \(a\in\{0,1\}\). Let \(\mu_x\) be the product law under which the outcome and factorial counts are independent with
\[
S_{ax}\sim\operatorname{Poi}(u q_{ax}),
\qquad
K_{ax}\sim\operatorname{Poi}(t s_{ax}),
\qquad a\in\{0,1\},
\]
and define the light polynomial cell contribution
\[
P_x
=
\frac{S_{1x}}{u}\widehat W_{1,L,x}
-
\frac{S_{0x}}{u}\widehat W_{0,L,x},
\]
where the factorial lift is taken at degree \(L\), scale \(B\), and intensity \(t\). Then
\[
\sum_{x:\,p_x>B}\pi_x\,\mathbb E_{\mu_x}\!\left[P_x^2\right]
\le
132\,n^{-12}.
\]
\end{lemmav}

\begin{lemmav}{lem:hybrid-heavy-branch-mean-bound}[Heavy branch mean bound]
Let \(d\) be a finite alphabet size, let \(\epsilon>0\), and let \(P\in\mathcal M_{d,\epsilon}\) as in \cref{obj:def:model-class}.  Let \(n\ge1\) and \(m\ge0\) be integers.  Assume:
\begin{itemize}
\item \textbf{(Calibration.)} The hybrid-estimator calibration predicate \(\mathsf C_\epsilon(n,m)\) in \cref{obj:def:mixed-estimator-handle} holds.
\item \textbf{(Block intensities.)} With \(M_0,n_p,n_f,m_p,m_f\) as in \cref{obj:def:mixed-estimator-handle}, set
\[
u=\frac{M_0}{8},\qquad
t_p=\frac{n_p+m_p}{8},\qquad
t=\frac{n_f+m_f}{8}.
\]
\item \textbf{(Cell laws.)} For each \(x\in[d]\), let \(\mu_x\) be the product law under which the outcome-block success counts \(S_{ax}\) are independent \(\operatorname{Poi}(u q_{ax})\) variables and the factorial-pool counts \(K_{ax}\) are independent \(\operatorname{Poi}(t s_{ax})\) variables, for \(a\in\{0,1\}\).
\end{itemize}
Let \(L=L(n)\), \(B=B(n,m,\epsilon)\), and \(k_0=k_0(n,m,\epsilon)\) be the degree, light-cell scale, and pilot threshold from \cref{obj:def:mixed-estimator-handle}.  For each \(x\in[d]\), let
\[
\pi_x=\mathbb P\{J_x\le k_0\},
\qquad
\eta_x=\mathbb P\{J_x>k_0\},
\]
where \(J_x=J_{0x}+J_{1x}\) has independent components \(J_{ax}\sim\operatorname{Poi}(t_p s_{ax})\).  Let \(P_x\) be the light polynomial cell contribution and \(H_x\) the heavy inverse-count cell contribution from \cref{obj:def:mixed-estimator-handle}, evaluated under \(\mu_x\).  Then
\[
\sum_{x\in[d]:\,p_x>B}
\pi_x\eta_x
\left(\mathbb E_{\mu_x}P_x-\mathbb E_{\mu_x}H_x\right)^2
\le 282\,n^{-12}.
\]
\end{lemmav}

\begin{lemmav}{lem:calibrated-hybrid-variance-bound}[Hybrid variance bound]
Let \(\epsilon\), the overlap parameter, satisfy \(0<\epsilon<1/2\). There is a constant \(C_{\mathrm{var}}>0\), depending only on \(\epsilon\), such that the following holds. Let \(n,m,d\) be integers and let \(P\) be an observed-data law on \([d]\times\{0,1\}^2\). Assume:
\begin{itemize}
\item \textbf{(Sample size.)} \(n\ge 1\).
\item \textbf{(Model class.)} \(P\in\mathcal M_{d,\epsilon}\), as in \cref{obj:def:model-class}.
\item \textbf{(Calibration.)} The finite calibration predicate \(\mathsf C_\epsilon(n,m)\) in \cref{obj:def:mixed-estimator-handle} holds.
\end{itemize}
Let
\[
M_0=\lfloor n/3\rfloor,\qquad
n_p=\bigl\lfloor (n-M_0)/2\bigr\rfloor,\qquad
n_f=n-M_0-n_p,\qquad
m_p=\lfloor m/2\rfloor,\qquad
m_f=m-m_p,
\]
and set
\[
u=M_0/8,\qquad
t_p=(n_p+m_p)/8,\qquad
t=(n_f+m_f)/8.
\]
Let \(L=L(n)\), \(B=B(n,m,\epsilon)\), \(k_0=k_0(n,m,\epsilon)\), and \(A_\star=7056\) be the degree, scale, pilot threshold, and factorial-moment constant from \cref{obj:def:mixed-estimator-handle}. Under the independent Poisson count law in which, independently over \(x\in[d]\) and \(a\in\{0,1\}\),
\[
S_{ax}\sim\operatorname{Poi}(u q_{ax}),\qquad
J_{ax}\sim\operatorname{Poi}(t_p s_{ax}),\qquad
K_{ax}\sim\operatorname{Poi}(t s_{ax}),
\]
let
\[
\widetilde Z
=
\sum_{x\in[d]}
\left(
\mathbf 1\{J_{0x}+J_{1x}\le k_0\}P_x
+
\mathbf 1\{J_{0x}+J_{1x}>k_0\}H_x
\right),
\]
where \(P_x\) and \(H_x\) are the light-cell polynomial and heavy inverse-count cell contributions specified in \cref{obj:def:mixed-estimator-handle}. Then
\[
\operatorname{Var}(\widetilde Z)
\le
4A_\star^L
\left\{
\frac{\min\{1,dB\}}{u}+dB^2
\right\}
+
C_{\mathrm{var}}\left(u^{-1}+t^{-1}\right)
+
\frac{8dB^2}{\epsilon^2L^4}
+
\frac{4}{\epsilon t}
+
414\,n^{-12}.
\]
\end{lemmav}

\begin{lemmav}{lem:calibrated-hybrid-mean-bias}[Calibrated hybrid mean bias]
Fix an overlap constant \(0<\epsilon<1/2\).  For every \(n,m,d\in\mathbb N\) and every observed-data law \(P\) on \([d]\times\{0,1\}^2\), suppose that:
\begin{itemize}
\item \textbf{(Sample size.)} \(1\le n\).
\item \textbf{(Model class.)} \(P\in\mathcal M_{d,\epsilon}\), as in \cref{obj:def:model-class}.
\item \textbf{(Calibration.)} The finite calibration predicate \(\mathsf C_\epsilon(n,m)\) in \cref{obj:def:mixed-estimator-handle} holds.
\end{itemize}
Let \(M_0,n_p,n_f,m_p,m_f\) be the block sizes from \cref{obj:def:mixed-estimator-handle}, and set
\[
u=\frac{M_0}{8},\qquad
t_p=\frac{n_p+m_p}{8},\qquad
t=\frac{n_f+m_f}{8}.
\]
Let \(L=L(n)\), \(B=B(n,m,\epsilon)\), and \(k_0=k_0(n,m,\epsilon)\) be the degree, light-cell scale, and pilot threshold from \cref{obj:def:mixed-estimator-handle}.  For each \(x\in[d]\), let the hybrid selected cell mean be the expectation of the cell contribution selected at threshold \(k_0\), degree \(L\), scale \(B\), and intensities \((u,t_p,t)\) under the independent Poisson count law in \cref{obj:def:mixed-estimator-handle}.  Then
\[
\left|
\sum_{x\in[d]}
\mathbb E_P[Z_x]
-\tau(P)
\right|
\le
\frac{2dB}{\epsilon L^2}+6n^{-12},
\]
where \(\tau(P)\) is the average treatment effect functional in \cref{obj:def:ate-functional}.
\end{lemmav}

\begin{definitionv}{synth_9}[Calibrated Poisson hybrid statistic]
Use the split sizes, intensities, constants, dyadic overlap certificate, and calibration predicate \(\mathsf C_\epsilon(n,m)\) from \(\cref{obj:def:mixed-estimator-handle}\). In the independent Poisson experiment with outcome-prefix length \(D_0\sim\operatorname{Poi}(u)\), pilot-prefix length \(D_p\sim\operatorname{Poi}(t_p)\), and factorial-prefix length \(D_f\sim\operatorname{Poi}(t)\), define \(Z^{\mathrm{mix}}\) as follows. If \(\mathsf C_\epsilon(n,m)\) is false, set
\[
Z^{\mathrm{mix}}:=0 .
\]
If \(\mathsf C_\epsilon(n,m)\) is true, return zero on the event that any Poisson prefix exceeds its corresponding pool. On the complementary valid-prefix event, let \(J_x\) be the pilot count in cell \(x\), let \(K_{ax}\) be the factorial-pool count in arm \(a\) and cell \(x\), and let \(S_{ax}\) be the outcome-block success count in arm \(a\) and cell \(x\). With
\[
H_L(z)=\frac{1-T_L(1-2z)}{2L^2},\qquad
E_L(z)=\frac{H_L(z)}{z},\qquad
G_L(z)=\frac{1-E_L(z)}{z},
\]
under polynomial continuation at zero, expand
\[
W_{a,L}(s_0,s_1)=\frac{s_0+s_1}{B}G_L(s_a/B)
=\sum_{r_0,r_1}c_{a,r_0,r_1}s_0^{r_0}s_1^{r_1}.
\]
For falling factorials \((k)_r\), put
\[
\widehat W_{a,L,x}
=\sum_{r_0,r_1}c_{a,r_0,r_1}
\frac{(K_{0x})_{r_0}(K_{1x})_{r_1}}{t^{r_0+r_1}},
\qquad
\widehat W^H_{a,x}
=\frac{K_{0x}+K_{1x}+1}{K_{ax}+1},
\]
and define the cell contribution
\[
Z_x=
\mathbf1\{J_x\le k_0\}
\left(\frac{S_{1x}}u\widehat W_{1,L,x}
-\frac{S_{0x}}u\widehat W_{0,L,x}\right)
+\mathbf1\{J_x>k_0\}
\left(\frac{S_{1x}}u\widehat W^H_{1,x}
-\frac{S_{0x}}u\widehat W^H_{0,x}\right).
\]
The calibrated Poisson hybrid statistic is the clipped aggregate
\[
Z^{\mathrm{mix}}
:=
\operatorname{clip}_{[-1,1]}\!\left(\sum_{x=1}^d Z_x\right)
\]
on the valid-prefix event, together with the preceding zero values on invalid prefixes or failed calibration.
\end{definitionv}

\begin{lemmav}{lem:poisson-hybrid-risk}[Hybrid Poisson risk]
Fix an overlap constant \(\epsilon\) with \(0<\epsilon<1/2\).  There is a constant \(C_\epsilon\in(0,\infty)\), depending only on \(\epsilon\), such that the following holds for every \(n,m,d\) and every observed-data law \(P\):
\begin{itemize}
\item \textbf{(Sample size.)} \(n\ge 1\).
\item \textbf{(Alphabet size.)} \(d\ge 2\).
\item \textbf{(Model class.)} \(P\in\mathcal M_{d,\epsilon}\), as in \cref{obj:def:model-class}.
\end{itemize}
Let \(M_0=\lfloor n/3\rfloor\), \(n_p=\lfloor(n-M_0)/2\rfloor\), \(n_f=n-M_0-n_p\), \(m_p=\lfloor m/2\rfloor\), and \(m_f=m-m_p\), and set
\[
u=\frac{M_0}{8},
\qquad
t_p=\frac{n_p+m_p}{8},
\qquad
t=\frac{n_f+m_f}{8}.
\]
In the independent Poisson count experiment with these intensities, let \(Z^{\mathrm{mix}}\) be the calibrated hybrid statistic when the finite calibration predicate \(\mathsf C_\epsilon(n,m)\) holds and the constant \(0\) otherwise.  Then its mean-squared error satisfies
\[
\mathbb E_P\!\left[\bigl(Z^{\mathrm{mix}}-\tau(P)\bigr)^2\right]
\le
C_\epsilon\, r_\epsilon(n,m,d),
\]
where \(r_\epsilon(n,m,d)\) is the frontier rate in \cref{obj:def:frontier-rate-handle}, equivalently
\[
r_\epsilon(n,m,d)
=
\min\!\left\{
1,\,
\frac1n+
\frac{d^2}{(n+m)^2\log^2(en)}
\right\}.
\]
\end{lemmav}

\begin{definitionv}{synth_6}[Three-pool Poisson count law]
Fix an observed-data law \(P\), intensities \(u,t_p,t\ge0\), and write
\[
q_{ax}=P(X=x,A=a,Y=1),
\qquad
s_{ax}=P(X=x,A=a).
\]
The law \(\mathbb P^{\mathrm{Pois}}_{P,u,t_p,t}\) is the joint distribution of independent count triples
\[
(S_{ax},J_{ax},K_{ax})_{a\in\{0,1\},\,x\in[d]},
\]
where
\[
S_{ax}\sim\operatorname{Poi}(u q_{ax}),\qquad
J_{ax}\sim\operatorname{Poi}(t_p s_{ax}),\qquad
K_{ax}\sim\operatorname{Poi}(t s_{ax}).
\]
Here \(S_{ax}\) is the outcome-success count from the outcome pool, \(J_{ax}\) is the pilot count, and \(K_{ax}\) is the factorial-pool treatment-covariate count. The product law makes all these counts independent across cells, arms, and pools.
\end{definitionv}

\begin{lemmav}{lem:fixed-sample-hybrid-transfer}[Fixed-sample hybrid transfer]
There is a universal constant \(C>0\) such that the following holds.  Let
\begin{itemize}
\item \textbf{(Overlap level.)} The overlap level \(\epsilon\) satisfies \(0<\epsilon<1/2\).
\item \textbf{(Sample sizes and alphabet.)} \(n\ge 1\), \(m\ge 0\), and \(d\ge 2\).
\item \textbf{(Model class.)} \(P\in\mathcal M_{d,\epsilon}\), with \(\mathcal M_{d,\epsilon}\) as in \cref{obj:def:model-class}.
\end{itemize}
For the mixed estimator \(\widehat\tau^{\mathrm{mix}}_{n,m,d}\) in \cref{obj:def:mixed-estimator-handle}, the estimator is measurable as a function of the annotation sample and
\[
\mathbb E_{P^{\otimes n}\otimes P_{XA}^{\otimes m}}
\left[
  \left\{\widehat\tau^{\mathrm{mix}}_{n,m,d}-\tau(P)\right\}^2
\right]
\le
2\,\mathcal R^{\mathrm{Pois}}_{\mathrm{hyb}}(P;n,m,\epsilon)
+\frac{C}{n}.
\]
Here \(\tau(P)\) is the average treatment effect functional in \cref{obj:def:ate-functional}.  Writing
\(M_0=\lfloor n/3\rfloor\), \(n_p=\lfloor(n-M_0)/2\rfloor\),
\(n_f=n-M_0-n_p\), \(m_p=\lfloor m/2\rfloor\), \(m_f=m-m_p\),
\(M_p=n_p+m_p\), \(M_f=n_f+m_f\), \(u=M_0/8\), \(t_p=M_p/8\), and
\(t=M_f/8\), the Poisson hybrid risk is
\[
\mathcal R^{\mathrm{Pois}}_{\mathrm{hyb}}(P;n,m,\epsilon)
=
\begin{cases}
\displaystyle
\int
\left\{
  Z^{\mathrm{mix}}(K)-\tau(P)
\right\}^2
\,d\mathbb P^{\mathrm{Pois}}_{P,u,t_p,t}(K),
&
\mathsf C_\epsilon(n,m),\\[1.2em]
\displaystyle
\tau(P)^2,
&
\text{otherwise,}
\end{cases}
\]
where \(Z^{\mathrm{mix}}\) is the calibrated hybrid statistic of \cref{obj:synth_9}.
\end{lemmav}

\begin{lemmav}{lem:shifted-chebyshev-quotient-coefficient-bound}[Shifted quotient coefficient bound]
Let \(L\) be a positive integer, and let \(G_L\) be the polynomial continuation at zero of \((1-E_L(z))/z\) from \cref{obj:def:mixed-estimator-handle}.  For a real polynomial \(p(z)=\sum_i a_i z^i\), write
\[
\|p\|_1=\sum_i |a_i|
\]
for its coefficient \(\ell^1\)-norm.  Then
\[
\|G_L\|_1\le 6^L .
\]
\end{lemmav}

\begin{lemmav}{lem:explicit-chebyshev-calibration}[Explicit Chebyshev calibration]
Let \(T_j\) be the Chebyshev polynomial of degree \(j\), set \(Q_j(z)=T_j(1-2z)\), and let \(\|\cdot\|_1\) denote coefficient \(\ell^1\)-norm.  For every integer \(L\ge2\), the polynomial continuations \(E_L\) and \(G_L\) satisfy the following calibration.

\begin{itemize}
\item \textbf{(Coefficient bounds.)}
\[
\|Q_L\|_1\le 7^L,
\qquad
2\|G_L\|_1\le 14^L .
\]
\item \textbf{(Factorial-lift second moment.)}
For every \(a\in\{0,1\}\), every \(t>0\), \(B>0\), and every \(s_0,s_1\ge0\) with \(L\le tB\), set \(p=s_0+s_1\).  If \(K_b\sim\operatorname{Poi}(ts_b)\) independently and \(\widehat W_{a,L}\) denotes the factorial lift of the light-cell polynomial based on \(G_L\), then
\[
E\widehat W_{a,L}^{2}
\le
7056^L\left(1+\frac{p}{B}\right)^{2L},
\]
and, whenever \(p\le B\),
\[
E\widehat W_{a,L}^{2}\le 7056^L .
\]
\item \textbf{(Interval identities.)}
For every \(z\in[0,1]\),
\[
0\le E_L(z),
\qquad
E_L(z)\le 1,
\qquad
zG_L(z)+E_L(z)=1 .
\]
Moreover, for every \(z\in(0,1]\),
\[
E_L(z)\le \frac{1}{L^2z}.
\]
\end{itemize}
\end{lemmav}

\begin{definitionv}{def:finite-small-instance}[Finite small instance \(\mathfrak C_{\mathrm{small}}\)]
Let
\[
\epsilon=\frac14,\qquad d=2,\qquad p_1=p_2=\frac12,\qquad e_1=e_2=\frac12,\qquad \mu_{01}=\mu_{02}=0.
\]
For \(n=m=1\), independence of the samples gives
\[
\mathbb E\!\left[
2\mathbf 1\{\mathcal L_1=((1,1,1))\}
\left(\mathbf 1\{\mathcal U_1=((1,0))\}+\mathbf 1\{\mathcal U_1=((1,1))\}\right)
\right]
=2q_{11}p_1,
\]
where \(q_{11}\) is the outcome-marked mass for treatment-covariate cell \((1,1)\). At \(\mu_{11}=1\), this equals \(1/4\), while the treated-cell target is
\[
p_1\mu_{11}=\frac12.
\]
Define \(\Pi_0^{\mathrm{small}}\) to put unit mass on the law with
\[
\mu_{11}=\mu_{12}=0,
\]
and define \(\Pi_1^{\mathrm{small}}\) to put unit mass on the law with
\[
\mu_{11}=1,\qquad \mu_{12}=0.
\]
Under both priors,
\[
P_{XA}(x,a)=\frac14
\quad\text{for all }(x,a),
\]
whereas
\[
\tau(P)=0
\]
under \(\Pi_0^{\mathrm{small}}\), and
\[
\tau(P)=\frac12
\]
under \(\Pi_1^{\mathrm{small}}\). Consequently, the auxiliary laws agree for every \(m\), and the total variation distance between the two labeled laws is
\[
1-\left(\frac34\right)^n.
\]
These identities constitute \(\mathfrak C_{\mathrm{small}}\).
\end{definitionv}

\begin{lemmav}{lem:finite-alternation-duality}[Finite alternation duality]
Fix an overlap constant \(\epsilon\) with \(0<\epsilon<1/2\), and set
\[
\kappa=\frac{1-2\epsilon}{\epsilon}.
\]
There exist constants \(\gamma_\epsilon>0\) and \(c_\epsilon>0\) such that, for every integer \(L\ge 2\) and every \(B>0\), with
\[
a=\frac{\gamma_\epsilon B}{L^2},
\]
one has \(a/\kappa\le B\), and there are distinct points \(p_0,\ldots,p_{L+1}\in [a/\kappa,B]\) and real weights \(w_0,\ldots,w_{L+1}\) satisfying
\[
\sum_{i=0}^{L+1}|w_i|=1,\qquad
\sum_{i=0}^{L+1}w_i p_i^j=0\quad\text{for every }0\le j\le L,
\]
and
\[
c_\epsilon\le \sum_{i=0}^{L+1}w_i\,\frac{p_i}{p_i+a}.
\]
\end{lemmav}

\begin{definitionv}{synth_20}[Common-marginal rare-cell recipe]
A common-marginal rare-cell recipe with \(k\) rare cells at scale \((a,B,L)\) is a finite random construction \(\mathcal R\) consisting of a generator law \(\operatorname{generator}_{\mathcal R}\) for a latent draw \(w\), together with two branch laws indexed by \(b\in\{0,1\}\). For each generated \(w\), the recipe specifies rare-cell raw masses \(p_i(w)\in[0,B]\), \(i=1,\ldots,k\), branch-specific outcome regressions, and any reservoir mass needed to normalize the construction. The two branches have the same treatment--covariate raw masses and treatment propensities, so after normalization they induce the same \(P_{XA}\) table, while their outcome regressions may differ.

For each branch \(b\in\{0,1\}\), let \(P_b(w)\) be the normalized law on \([d]\times\{0,1\}^2\) induced by the recipe, and let
\[
\operatorname{rawTarget}_{\mathcal R}(b,w)
\]
be the unnormalized target contribution assigned by the rare cells in branch \(b\). The branch center is
\[
\theta_b=\int \operatorname{rawTarget}_{\mathcal R}(b,w)\,
d\operatorname{generator}_{\mathcal R}(w),
\]
and the raw-center separation is \(\Delta=|\theta_1-\theta_0|\). The recipe induces, for each branch \(b\), the mixture experiment obtained by first drawing
\(w\sim\operatorname{generator}_{\mathcal R}\), then sampling from \(P_b(w)\) for labeled observations and from the common marginal \(P_{b,XA}(w)\) for auxiliary observations. Its raw Poisson version uses labeled intensity \(u\) and auxiliary intensity \(v\) with the same generated law \(P_b(w)\), so the treatment--covariate component is common across the two branches and the likelihood discrepancy is carried by the labeled outcome-bearing component.
\end{definitionv}

\begin{lemmav}{lem:common-marginal-recipe-target-concentration}[Prior target concentration]
Fix \(\epsilon\in\mathbb R\), a common-marginal calibration at overlap level \(\epsilon\) with dual-gap constant \(c_\epsilon\), and parameters \(n,m,d,L,k\in\mathbb N\) and \(B,a,u,v,C,\rho\in\mathbb R\).  Let \(\mathcal R\) be a common-marginal rare-cell recipe with atom-generator measure \(\Gamma\), branch laws \(P_b(w)\), raw targets \(r_b(w)\), raw total mass \(S(w)\), and branch priors
\[
\Pi_b=\Gamma\circ P_b^{-1},
\qquad
\theta_b=\int r_b(w)\,d\Gamma(w),
\qquad b\in\{0,1\}.
\]
Assume the recipe bounds hold at intensities \((u,v)\) with constants \((C,\rho)\) and dual gap \(c_\epsilon\):
\begin{itemize}
\item \textbf{(Finite-atom construction.)}  The recipe uses the finite-atom construction.
\item \textbf{(Model support.)}  Each prior is supported on \(\mathcal M_{d,\epsilon}\) of \cref{obj:def:model-class}.
\item \textbf{(Common table.)}  The two priors induce the same law of the random treatment-covariate table.
\item \textbf{(Separation.)}  The raw centers satisfy
\[
c_\epsilon k a\le |\theta_1-\theta_0|.
\]
\item \textbf{(Mass and variance bounds.)}  Under the usual measurability and integrability conditions,
\[
\int S(w)\,d\Gamma(w)=1,\qquad
\operatorname{Var}_\Gamma(S)\le CkBa,\qquad
\max_{b\in\{0,1\}}\operatorname{Var}_\Gamma(r_b)\le CkBa.
\]
\item \textbf{(Mixture proximity.)}  The two raw mixture laws at labeled intensity \(u\) and auxiliary intensity \(v\) have total variation distance at most
\[
Cuka\rho^L.
\]
\end{itemize}
Then for every branch \(b\in\{0,1\}\) and every \(\Delta\) with \(0<\Delta\le2\), the branch prior concentrates its average treatment effect \(\tau(P)\), defined in \cref{obj:def:ate-functional}, around its raw center:
\[
\Pi_b\!\left\{P:\frac{\Delta}{4}<|\tau(P)-\theta_b|\right\}
\le
\frac{512\,CkBa}{\Delta^2}.
\]
\end{lemmav}

\begin{definitionv}{synth_21}[Raw Poisson prior-predictive mixture laws]
Fix a common-marginal rare-cell recipe \(\mathcal R\) with atom generator \(\Gamma\), branch laws \(P_b(w)\) for \(b\in\{0,1\}\), and branch priors \(\Pi_b=\Gamma\circ P_b^{-1}\). At labeled and auxiliary Poisson intensities \(u\) and \(v\), define \(\mathsf Q_b\) to be the branch-\(b\) raw prior-predictive law obtained by first drawing \(w\sim\Gamma\), forming the raw branch law \(P_b(w)\), and then sampling the independent raw Poisson experiment with labeled intensity \(u\) and auxiliary intensity \(v\) from that raw law and its treatment--covariate marginal. Equivalently,
\[
\mathsf Q_b
=
\int
\Bigl\{
\operatorname{PoiExp}_{u,v}^{\mathrm{raw}}\bigl(P_b(w)\bigr)
\Bigr\}\,d\Gamma(w)
=
\int
\operatorname{PoiExp}_{u,v}^{\mathrm{raw}}(P)\,d\Pi_b(P),
\qquad b\in\{0,1\},
\]
where \(\operatorname{PoiExp}_{u,v}^{\mathrm{raw}}(P)\) denotes the product Poisson law for the outcome-labeled raw cell counts at intensity \(u\) and the auxiliary treatment--covariate raw cell counts at intensity \(v\).
\end{definitionv}

\begin{lemmav}{lem:common-marginal-recipe-scale-tail}[Small-scale prior tail]
Fix \(\epsilon\in\mathbb R\), a common-marginal calibration with positive constants including dual-gap constant \(c_\epsilon\le 1/4\), integers \(n,m,d,L,k\ge0\), and reals \(B,a,u,v,C,\rho\).  Let \(\mathcal R\) be a common-marginal rare-cell recipe with atom-generator \(\Gamma\), branch laws \(P_b(w)\), branch priors \(\Pi_b=\Gamma\circ P_b^{-1}\), raw total mass \(S(w)\), raw targets \(r_b(w)\), and measurable raw scale \(s_{\mathcal R}\) satisfying \(s_{\mathcal R}(P_b(w))=S(w)\) for \(\Gamma\)-almost every \(w\).

Assume that the recipe bounds hold for \(\mathcal R\) at \((u,v,c_\epsilon,C,\rho)\):
\begin{itemize}
\item \textbf{(Finite-atom construction.)} The recipe uses the finite-atom construction.
\item \textbf{(Model support and common table.)} The two priors are supported on \(\mathcal M_{d,\epsilon}\) from \cref{obj:def:model-class}, and they induce the same distribution of the random treatment-covariate table \(P_{XA}\).
\item \textbf{(Center separation.)} If \(\theta_b\) denotes the raw prior center for branch \(b\), then
\[
c_\epsilon k a \le |\theta_1-\theta_0|.
\]
\item \textbf{(Mass and variance bounds.)} Under the usual measurability and integrability conditions,
\[
\int S\,d\Gamma=1,\qquad
\operatorname{Var}_{\Gamma}(S)\le CkBa,\qquad
\max\{\operatorname{Var}_{\Gamma}(r_0),\operatorname{Var}_{\Gamma}(r_1)\}\le CkBa.
\]
\item \textbf{(Mixture distance.)} The two raw Poisson mixture laws generated by the recipe satisfy
\[
\operatorname{TV}(\mathsf Q_0,\mathsf Q_1)\le Cuka\rho^L.
\]
\end{itemize}
Then, for each branch \(b\in\{0,1\}\),
\[
\Pi_b\!\left\{P:\max\{s_{\mathcal R}(P),0\}<\frac12\right\}
\le 4CkBa .
\]
\end{lemmav}

\begin{lemmav}{lem:common-marginal-recipe-transfer}[Recipe transfer bound]
Let \(0<\epsilon<1/2\), where \(\epsilon\) is the overlap level.  Fix a common-marginal calibration at level \(\epsilon\), consisting of positive constants including \(\gamma_\epsilon\) and dual gap \(\delta_\epsilon\le 1/4\), a constant \(C_\epsilon>0\), and a real number \(\rho_\epsilon\).  Let \(n,m,d,L,k\) be nonnegative integers, let \(B,a\in\mathbb R\), and let \(\mathcal R\) be a common-marginal rare-cell recipe with \(k\) rare cells at scale \((a,B,L)\).  For each branch \(b\in\{0,1\}\), define the raw prior center
\[
\theta_b=\int \operatorname{rawTarget}_{\mathcal R}(b,w)\,dG_{\mathcal R}(w),
\]
where \(G_{\mathcal R}\) is the atom-generator measure of the recipe, and put
\[
\Delta=|\theta_1-\theta_0|.
\]
Assume:
\begin{itemize}
\item \textbf{(Sample and dimension.)} \(n\ge 1\) and \(d\ge 2\).
\item \textbf{(Poisson-prior data.)} With
\[
u=65536\,n
\qquad\text{and}\qquad
v=65536\,(n+m),
\]
the recipe has bandwidth \(B>0\), shift
\[
a=\frac{\gamma_\epsilon B}{L^2},
\]
rare-shift size
\[
ka\le \delta_\epsilon,
\]
and the recipe intensities agree with \(u\) and \(v\).
\item \textbf{(Bandwidth and likelihood bounds.)} The realized scale satisfies
\[
(u+v)B\le b_\epsilon L,
\]
where \(b_\epsilon\) is the calibration's uniform-bandwidth constant, and the recipe satisfies the common-marginal Poisson-prior bounds at gap \(\delta_\epsilon\), intensity pair \((u,v)\), variance constant \(C_\epsilon\), and decay constant \(\rho_\epsilon\).
\item \textbf{(Mixture affinity.)} The two prior-predictive raw Poisson-count laws generated by the two recipe priors and the recipe's experiment kernel have total variation distance at most
\[
\frac18 .
\]
\item \textbf{(Variance scale.)} The recipe scale obeys
\[
kBa\le \frac{\Delta^2}{65536\,C_\epsilon}.
\]
\end{itemize}
Then the minimax risk of \cref{obj:def:minimax-risk} satisfies
\[
R_\epsilon(n,m,d)\ge \frac{\Delta^2}{65536}.
\]
\end{lemmav}

\begin{lemmav}{lem:fixed-sample-common-marginal-transfer}[Fixed-sample transfer]
There exist numerical constants \(C_{\mathrm{fixed}}>0\) and \(c_{\mathrm{risk}}>0\) such that the following holds.  For every \(0<\epsilon<1/2\) and every \(C_\epsilon>0\), there is a constant \(c_{\epsilon,C_\epsilon}>0\) such that, for every common-marginal calibration at overlap level \(\epsilon\) with positive calibration constants and dual gap \(\delta_\epsilon\le 1/4\), every \(\rho\in\mathbb R\), every \(n,m,d,L,k\in\mathbb N\), every \(B,a\in\mathbb R\), and every common-marginal recipe \(R\) at these indices, the implication below holds.

Write \(G_R\) for the recipe's rare-cell generator measure and \(\theta_{R,b}(w)\) for its raw target on branch \(b\).  Let
\[
\Delta
=
\int \theta_{R,1}(w)\,dG_R(w)
-
\int \theta_{R,0}(w)\,dG_R(w).
\]
Assume:
\begin{itemize}
\item \textbf{(Sample and dimension.)} \(n\ge 1\) and \(d\ge 2\).
\item \textbf{(Witness data.)} The recipe carries a common-marginal Poisson-prior witness at labeled and auxiliary intensities
\[
u=C_{\mathrm{fixed}}n,
\qquad
v=C_{\mathrm{fixed}}(n+m).
\]
Equivalently, the witness supplies
\[
(u+v)B\le b_\epsilon L,
\qquad
B>0,
\qquad
a=\gamma_\epsilon B/L^2,
\qquad
ka\le \delta_\epsilon,
\]
identifies the recipe's labeled and auxiliary intensities with \(u\) and \(v\), and satisfies the corresponding common-marginal bounds at dual gap \(\delta_\epsilon\), variance-bound constant \(C_\epsilon\), and likelihood-decay parameter \(\rho\).
\item \textbf{(Testing distance.)} The raw prior-predictive mixture laws \(M_0\) and \(M_1\) generated by the two branches of the recipe satisfy
\[
\operatorname{TV}(M_0,M_1)\le \frac18 .
\]
\item \textbf{(Concentration scale.)} The realized scale obeys
\[
kBa\le c_{\epsilon,C_\epsilon}\Delta^2 .
\]
\end{itemize}
Then the minimax risk \(R_\epsilon(n,m,d)\) from \cref{obj:def:minimax-risk} satisfies
\[
c_{\mathrm{risk}}\Delta^2
\le
R_\epsilon(n,m,d).
\]
\end{lemmav}

\begin{lemmav}{lem:parametric-label-floor}[Parametric label floor]
Let \(R_\epsilon(n,m,d)\) be the minimax risk in \cref{obj:def:minimax-risk}. There exists a universal constant \(c>0\) such that, whenever
\begin{itemize}
\item \textbf{(Overlap level.)} \(0<\epsilon<1/2\),
\item \textbf{(Labeled sample size.)} \(n\ge 1\),
\item \textbf{(Auxiliary sample size.)} \(m\) is a nonnegative integer,
\item \textbf{(Alphabet size.)} \(d\ge 2\),
\end{itemize}
the minimax risk satisfies
\[
R_\epsilon(n,m,d)\ge \frac{c}{n}.
\]
\end{lemmav}

\begin{lemmav}{lem:finite-product-padding-identity}[Finite product padding identity]
Let \(I\) and \(J\) be finite index sets, and let \(\mathsf Y\) be a measurable space whose singletons are measurable. Suppose the following conditions hold:
\begin{itemize}
\item \textbf{(Embedding.)} \(f:I\to J\) is injective.
\item \textbf{(Distinguished coordinate.)} \(r\in J\) satisfies \(r\ne f(i)\) for every \(i\in I\).
\item \textbf{(Coordinate laws.)} \(\mu_i\), \(i\in I\), and \(\nu\) are probability measures on \(\mathsf Y\), and \(y_0\in\mathsf Y\).
\end{itemize}
Define \(\mathsf A:\mathsf Y^I\times\mathsf Y\to\mathsf Y^J\) by
\[
\mathsf A(z,y)(j)=
\begin{cases}
y, & j=r,\\
z_i, & j=f(i)\ \text{for the unique } i\in I,\\
y_0, & j\in J\setminus\bigl(\{r\}\cup f(I)\bigr).
\end{cases}
\]
Define the coordinate laws \(\mathsf H_j\), \(j\in J\), by
\[
\mathsf H_j=
\begin{cases}
\nu, & j=r,\\
\mu_i, & j=f(i)\ \text{for the unique } i\in I,\\
\delta_{y_0}, & j\in J\setminus\bigl(\{r\}\cup f(I)\bigr).
\end{cases}
\]
Then
\[
\mathsf A_\#\left(\left(\bigotimes_{i\in I}\mu_i\right)\otimes\nu\right)
=
\bigotimes_{j\in J}\mathsf H_j .
\]
\end{lemmav}

\begin{lemmav}{lem:common-marginal-uniform-intensity}[Uniform intensity common marginal]
For every overlap level \(0<\epsilon<1/2\), there are positive constants
\[
r_\epsilon,\quad D_\epsilon,\quad b_{0,\epsilon},\quad b_\epsilon,\quad
\gamma_\epsilon,\quad q_\epsilon,\quad c_\epsilon,\quad \delta_\epsilon,\quad C_\epsilon
\]
and a number \(0<\rho_\epsilon<1\), with \(c_\epsilon\le 1/4\), such that the following calibration facts hold:
\begin{itemize}
\item \textbf{(Bandwidth slack.)} \(131072\,b_{0,\epsilon}\le b_\epsilon\).
\item \textbf{(Degree floor.)} \(2\le D_\epsilon\).
\item \textbf{(Rare-cell scale.)}
\[
q_\epsilon\gamma_\epsilon b_{0,\epsilon}\le c_\epsilon,
\qquad
q_\epsilon=
\frac{c_\epsilon}{2\gamma_\epsilon b_{0,\epsilon}},
\qquad
4\le q_\epsilon .
\]
\item \textbf{(Dimension scale.)}
\[
\delta_\epsilon r_\epsilon\le \frac1{65536},
\qquad
\delta_\epsilon\le \frac{c_\epsilon^4}{16},
\qquad
\delta_\epsilon\le
\left(\frac{c_\epsilon\gamma_\epsilon b_{0,\epsilon}}
{2(D_\epsilon+1)}\right)^2 .
\]
\item \textbf{(Variance and counting calibration.)}
\[
8(D_\epsilon+1)b_{0,\epsilon}
\le \frac{c_\epsilon^3}{65536\,C_\epsilon},
\]
\[
\frac{4(D_\epsilon+1)}
{(1/(65536\,C_\epsilon))c_\epsilon^2\gamma_\epsilon}
\le q_\epsilon,
\]
and
\[
\left\{
\frac{4(D_\epsilon+1)^2}
{(1/(65536\,C_\epsilon))c_\epsilon^2\gamma_\epsilon}
\right\}^2
\le r_\epsilon .
\]
\item \textbf{(Mixture decay.)} For every integer \(n\ge1\),
\[
C_\epsilon(65536n)c_\epsilon
\rho_\epsilon^{\lceil D_\epsilon\log(en)\rceil}
\le \frac18 .
\]
\end{itemize}

For every \(n,m,L,d,k\) and every \(u,v,B,a\), if
\begin{itemize}
\item \(n\ge1\), \(L\ge2\), and \(d\ge2\);
\item \(u\ge0\), \(v\ge0\), and \(u+v>0\);
\item \(B>0\), \((u+v)B\le b_\epsilon L\), and
\[
a=\frac{\gamma_\epsilon B}{L^2};
\]
\item \(ka\le c_\epsilon\) and \(k+1\le d\),
\end{itemize}
then there is a common-marginal rare-cell recipe with \(k\) rare cells at scale \((a,B,L)\) whose labeled and auxiliary Poisson intensities are \(u\) and \(v\). Its null and alternative priors are the two branch priors of the recipe, and they satisfy:
\begin{itemize}
\item the recipe is in the finite-atom construction regime, and both priors are supported on \(\mathcal M_{d,\epsilon}\);
\item the induced random treatment-covariate tables \(P_{XA}\) have exactly the same law under the two priors;
\item the raw total mass has mean one;
\item the raw target centers obey
\[
c_\epsilon ka
\le
\left|
\theta_1-\theta_0
\right|,
\]
where \(\theta_1\) and \(\theta_0\) are the alternative and null raw prior centers;
\item the raw total-mass variance is at most
\[
C_\epsilon kBa,
\]
and the larger of the two raw target variances is at most
\[
C_\epsilon kBa;
\]
\item in the raw Poisson experiment with labeled intensity \(u\) and auxiliary intensity \(v\), the two prior-predictive mixture laws have total variation distance at most
\[
C_\epsilon uka\,\rho_\epsilon^L .
\]
\end{itemize}
\end{lemmav}

\begin{lemmav}{lem:common-marginal-canonical-certificate-family}[Canonical certificate family]
There exist real constants \(C_{\mathrm{fixed}}\) and \(c_{\mathrm{risk}}\) with
\[
0<C_{\mathrm{fixed}},
\qquad
0<c_{\mathrm{risk}},
\]
such that, for every overlap level \(\epsilon\) with \(0<\epsilon<1/2\), there are a common-marginal calibration in the sense of \cref{obj:def:common-marginal-prior-handle} and real constants
\[
C_\epsilon,\rho_\epsilon,c_{\mathrm{small}},c_{\mathrm{floor}}
\]
with the following properties.  Write \(b_0\) for the calibration bandwidth constant, \(b_\epsilon\) for the uniform bandwidth constant, \(\gamma_\epsilon\) for the shift constant, \(c_\epsilon\) for the dual-gap constant, and \(C_d\) for the dimension constant supplied by the calibration.  These calibration constants are positive, and
\[
c_\epsilon\le \frac14 .
\]
The auxiliary constants satisfy
\[
0<C_\epsilon,
\qquad
\rho_\epsilon\in(0,1),
\qquad
0<c_{\mathrm{small}},
\qquad
0<c_{\mathrm{floor}},
\qquad
c_{\mathrm{floor}}\le c_{\mathrm{risk}} .
\]
Moreover:
\begin{itemize}
\item \textbf{(Fixed-sample transfer.)}
For every \(n,m,d\) with \(n\ge1\) and \(d\ge2\), and every common-marginal prior handle \(H\) at \((n,m,d)\) in the sense of \cref{obj:def:common-marginal-prior-handle}, let \(L,B,a,k\) denote the degree, bandwidth, shift, and rare-cell count components of \(H\), and let \(\theta_0(H)\) and \(\theta_1(H)\) denote the two raw prior centers of its recipe.  If the recipe of \(H\) carries the exact common-marginal Poisson-prior construction at bandwidth constant \(b_0\), information size \(n+m\), labeled intensity
\[
u=C_{\mathrm{fixed}}n,
\]
auxiliary intensity
\[
v=C_{\mathrm{fixed}}(n+m),
\]
and constants \((C_\epsilon,\rho_\epsilon)\), including the corresponding finite-support prior data, exact common treatment-covariate marginal table, moment-matching identities through degree \(L\), raw-center separation data, and Poisson likelihood-ratio bound at rate \(C_\epsilon u k a\rho_\epsilon^L\), if the two raw prior-predictive Poisson experiment laws supplied by that recipe have total variation distance at most \(1/8\), and if
\[
kBa
\le
c_{\mathrm{small}}
\left|\theta_1(H)-\theta_0(H)\right|^2,
\]
then
\[
c_{\mathrm{risk}}
\left|\theta_1(H)-\theta_0(H)\right|^2
\le
R_\epsilon(n,m,d),
\]
where \(R_\epsilon(n,m,d)\) is the minimax risk from \cref{obj:def:minimax-risk}.

\item \textbf{(Canonical handles.)}
For every \(n,m,d\) with \(n\ge1\) and \(d\ge2\), there exists a common-marginal prior handle \(H\) at \((n,m,d)\).  Let \(L,B,a,k\) denote its degree, bandwidth, shift, and rare-cell count components, set
\[
\Delta_H=c_\epsilon ka,
\qquad
D(n,m,d)=\min\left\{1,\frac{d^2}{(n+m)^2\log(en)^2}\right\},
\]
and let \(\theta_0(H)\) and \(\theta_1(H)\) denote the two raw prior centers of its recipe.  The same constants satisfy
\[
L\ge2,
\qquad
2C_{\mathrm{fixed}}b_0\le b_\epsilon,
\]
\[
\bigl(C_{\mathrm{fixed}}n+C_{\mathrm{fixed}}(n+m)\bigr)B
\le
b_\epsilon L,
\]
\[
\mathcal L(P_{XA}\mid \Pi_0(H))
=
\mathcal L(P_{XA}\mid \Pi_1(H)),
\qquad
\frac{c_{\mathrm{floor}}}{n}\le R_\epsilon(n,m,d).
\]
On the finite-atom branch, the handle also satisfies
\[
ka\le c_\epsilon,
\]
and its recipe carries the exact common-marginal Poisson-prior construction at bandwidth constant \(b_0\), information size \(n+m\), labeled intensity \(u=C_{\mathrm{fixed}}n\), auxiliary intensity \(v=C_{\mathrm{fixed}}(n+m)\), and constants \((C_\epsilon,\rho_\epsilon)\), including the finite-support prior data, exact common treatment-covariate marginal table, moment-matching identities through degree \(L\), raw-center separation data, and Poisson likelihood-ratio bound at rate \(C_\epsilon u k a\rho_\epsilon^L\).  The recipe bounds hold at these intensities: both priors are supported on \(\mathcal M_{d,\epsilon}\) from \cref{obj:def:model-class}, they induce the same random \(P_{XA}\) law, their raw-center separation is at least \(c_\epsilon ka\), the raw total mass of the generator has mean one, the raw total mass and raw targets are square-integrable with generator variances at most \(C_\epsilon kBa\), and the raw mixture total variation is at most \(C_\epsilon u k a\rho_\epsilon^L\).  Moreover, the two raw prior-predictive Poisson experiment laws supplied by the recipe have total variation distance at most \(1/8\), and
\[
\Delta_H\le
\left|\theta_1(H)-\theta_0(H)\right|,
\qquad
kBa\le c_{\mathrm{small}}\Delta_H^2,
\]
\[
c_{\mathrm{floor}}\Delta_H^2\le R_\epsilon(n,m,d),
\qquad
C_d\,D(n,m,d)\le \Delta_H^2.
\]
On the parametric branch,
\[
C_d\,D(n,m,d)\le \frac{c_{\mathrm{floor}}}{n}.
\]
\end{itemize}
\end{lemmav}

\begin{definitionv}{def:known-marginal-experiment}[Known-marginal experiment \(\mathcal E^{\mathrm{known}}_{n,d,\epsilon}\)]
For \(n\ge1\) and \(d\ge2\), the known-marginal experiment is
\[
\mathcal E^{\mathrm{known}}_{n,d,\epsilon}
:=
\{(P^{\otimes n},P_{XA}):P\in\mathcal M_{d,\epsilon}\}.
\]
The first coordinate generates \(\mathcal L_n\), and the second coordinate is the exact \(2d\)-cell table observed as deterministic side information.
\end{definitionv}

\begin{lemmav}{lem:finite-side-cover-approximate-decision}[Finite side cover]
Let \(\Theta\) be a nonempty compact topological space, let \(\mathsf Z\) be a finite set, let \(\mathsf C\) be a finite nonempty set, and let \(l\le u\). Suppose the following conditions hold.
\begin{itemize}
\item \textbf{(Laws.)} For each \(\theta\in\Theta\), \(\Lambda_\theta\) is a probability vector on \(\mathsf Z\) and \(\Xi_\theta\) is a probability vector on \(\mathsf C\).
\item \textbf{(Target.)} For each \(\theta\in\Theta\), \(\upsilon_\theta\in[l,u]\).
\item \textbf{(Continuity.)} For every \(z\in\mathsf Z\) and \(c\in\mathsf C\), the maps \(\theta\mapsto\Lambda_\theta(z)\), \(\theta\mapsto\Xi_\theta(c)\), and \(\theta\mapsto\upsilon_\theta\) are continuous.
\end{itemize}
For probability vectors \(\nu,\nu'\) on \(\mathsf C\), set
\[
D_\infty(\nu,\nu')=\max_{c\in\mathsf C}|\nu(c)-\nu'(c)|.
\]
Define
\[
\mathfrak R_\star
=
\inf_{\psi}\sup_{\theta\in\Theta}
\sum_{z\in\mathsf Z}\Lambda_\theta(z)
\bigl(\psi(z,\Xi_\theta)-\upsilon_\theta\bigr)^2,
\]
where \(\psi\) ranges over all maps
\[
\psi:\mathsf Z\times\mathsf{Prob}(\mathsf C)\to[l,u].
\]
Then for every \(\delta>0\), there exist an integer \(k\ge1\), centers
\(\theta^\circ_1,\ldots,\theta^\circ_k\in\Theta\), decisions
\(\beta_1,\ldots,\beta_k:\mathsf Z\to[l,u]\), radii
\(\operatorname{rad}_1,\ldots,\operatorname{rad}_k\), and a margin \(\eta>0\) such that
\[
\eta<\operatorname{rad}_j
\]
for every \(j\), every \(\theta\in\Theta\) admits some \(j\in\{1,\ldots,k\}\) with
\[
D_\infty(\Xi_\theta,\Xi_{\theta^\circ_j})<\operatorname{rad}_j-\eta,
\]
and, for every \(j\in\{1,\ldots,k\}\) and every \(\theta\in\Theta\),
\[
D_\infty(\Xi_\theta,\Xi_{\theta^\circ_j})<\operatorname{rad}_j
\quad\Longrightarrow\quad
\sum_{z\in\mathsf Z}\Lambda_\theta(z)
\bigl(\beta_j(z)-\upsilon_\theta\bigr)^2
\le
\mathfrak R_\star+\delta.
\]
\end{lemmav}

\begin{lemmav}{lem:finite-side-information-convergence}[Finite side-information convergence]
Let the overlap parameter \(\epsilon\) satisfy \(0<\epsilon<1/2\), let \(n\ge0\), and let \(d\ge2\).  Let
\[
\Theta_{\epsilon,d}
=
\Bigl\{(\pi,\varrho,\omega):
\pi\in\mathsf{Prob}([d]),\
\varrho_x\in[\epsilon,1-\epsilon],\
\omega_{\alpha x}\in[0,1]\ \text{for }x\in[d],\ \alpha\in\{0,1\}
\Bigr\}.
\]
For \(\theta=(\pi,\varrho,\omega)\in\Theta_{\epsilon,d}\), set
\[
P_\theta(x,\alpha,y)
=
\pi_x\varrho_x^\alpha(1-\varrho_x)^{1-\alpha}
\omega_{\alpha x}^y(1-\omega_{\alpha x})^{1-y},
\]
let \(\lambda_\theta=P_\theta^{\otimes n}\) be the labeled-sample law on \(([d]\times\{0,1\}^2)^n\), define
\[
\xi_\theta(x,\alpha)=\sum_{y\in\{0,1\}}P_\theta(x,\alpha,y),
\qquad
\upsilon_\theta=\sum_{x\in[d]}\pi_x(\omega_{1x}-\omega_{0x}).
\]
For each \(m\ge0\), define
\[
\mathfrak R_m
=
\inf_{\varphi:([d]\times\{0,1\}^2)^n\times([d]\times\{0,1\})^m\to[-1,1]}
\sup_{\theta\in\Theta_{\epsilon,d}}
\int
\bigl(\varphi(Z,U)-\upsilon_\theta\bigr)^2\,
d\bigl(\lambda_\theta\otimes\xi_\theta^{\otimes m}\bigr)(Z,U),
\]
and define
\[
\mathfrak R_\star
=
\inf_{\psi:([d]\times\{0,1\}^2)^n\times\mathsf{Prob}([d]\times\{0,1\})\to[-1,1]}
\sup_{\theta\in\Theta_{\epsilon,d}}
\int
\bigl(\psi(Z,\xi_\theta)-\upsilon_\theta\bigr)^2\,
d\lambda_\theta(Z).
\]
Then
\[
\lim_{m\to\infty}\mathfrak R_m=\mathfrak R_\star .
\]
\end{lemmav}

\begin{lemmav}{lem:finite-side-measurable-comparison}[Empirical side comparison]
Let \(\Theta\) be a nonempty parameter set, let \(\mathcal Z\) and \(\mathcal C\) be nonempty finite sets, and let \(\beta_-\le \beta_+\) be real numbers. Suppose that, for each \(\theta\in\Theta\), the following objects are given:
\begin{itemize}
\item \textbf{(Observed side law.)} A probability vector \(\lambda_\theta\) on \(\mathcal Z\).
\item \textbf{(Empirical side law.)} A probability vector \(\xi_\theta\) on \(\mathcal C\).
\item \textbf{(Target.)} A number \(\upsilon_\theta\in[\beta_-,\beta_+]\).
\end{itemize}
For each integer \(m\ge0\), define
\[
\mathfrak R_m
=
\inf_{\varphi:\mathcal Z\times\mathcal C^m\to[\beta_-,\beta_+]}
\sup_{\theta\in\Theta}
\sum_{z\in\mathcal Z}\lambda_\theta(z)
\sum_{v\in\mathcal C^m}
\left(\prod_{j=1}^m \xi_\theta(v_j)\right)
\bigl(\varphi(z,v)-\upsilon_\theta\bigr)^2 .
\]
Define also
\[
\widetilde{\mathfrak R}_\star
=
\inf_{\psi}
\sup_{\theta\in\Theta}
\sum_{z\in\mathcal Z}\lambda_\theta(z)
\bigl(\psi(z,\xi_\theta)-\upsilon_\theta\bigr)^2 ,
\]
where the infimum ranges over all maps
\[
\psi:\mathcal Z\times\mathbb R^{\mathcal C}\to[\beta_-,\beta_+]
\]
such that, for each \(z\in\mathcal Z\), the section \(w\mapsto\psi(z,w)\) is Borel measurable. Then, for every integer \(m\ge0\),
\[
\widetilde{\mathfrak R}_\star\le \mathfrak R_m .
\]
\end{lemmav}

\begin{lemmav}{lem:annotation-frontier-strict-improvement}[Strict frontier improvement]
Let \((m_n)_{n\ge 1}\) and \((d_n)_{n\ge 1}\) be sequences of nonnegative integers such that \(d_n\ge 2\) for every \(n\ge 1\). For each \(n\ge 1\), set \(N_n=n+m_n\) and \(\ell_n=\log(en)\), and let \(r_\epsilon(n,m,d)\) be the annotation-frontier rate in \cref{obj:def:frontier-rate-handle}. Then, as \(n\to\infty\),
\[
r_\epsilon(n,m_n,d_n)
=
o\!\left(r_\epsilon(n,0,d_n)\right)
\]
if and only if
\[
\frac{d_n}{\sqrt n\,\ell_n}\to\infty
\]
and
\[
\frac{d_n/(N_n\ell_n)}
{\min\{1,d_n/(n\ell_n)\}}
\to 0.
\]
\end{lemmav}
